% SPDX-License-Identifier: MIT-0
% Merged research draft. Incorporates material originally written as three
% separate packages; see Section 1.5 (Provenance) and Appendix E.
\documentclass[11pt]{article}
\usepackage[T1]{fontenc}
\usepackage[utf8]{inputenc}
\usepackage{lmodern}
\usepackage[a4paper,margin=25mm,headheight=14pt]{geometry}
\usepackage{amsmath,amssymb,amsthm,mathtools}
\usepackage{microtype,booktabs,array,longtable}
\usepackage[dvipsnames]{xcolor}
\usepackage{enumitem,fancyhdr}
\usepackage[colorlinks=true,linkcolor=MidnightBlue,citecolor=MidnightBlue,
 urlcolor=MidnightBlue,bookmarksnumbered=true]{hyperref}
\usepackage{listings}
\definecolor{Ink}{HTML}{19384B}
\definecolor{CodeBackground}{HTML}{F4F5EF}
\lstset{basicstyle=\small\ttfamily,backgroundcolor=\color{CodeBackground},
 breaklines=true,columns=fullflexible,keepspaces=true,showstringspaces=false,
 frame=single,rulecolor=\color{black!15},xleftmargin=4pt,xrightmargin=4pt}
\pagestyle{fancy}
\fancyhf{}
\fancyhead[L]{\small Support reciprocity for preorder posets}
\fancyhead[R]{\small Research draft}
\fancyfoot[C]{\thepage}
\renewcommand{\headrulewidth}{0.3pt}
\setlength{\parskip}{0.3em}
\setlength{\parindent}{1.2em}
\setlist{itemsep=0.2em,topsep=0.4em}
\emergencystretch=2em
\numberwithin{equation}{section}
\newtheorem{theorem}{Theorem}[section]
\newtheorem{lemma}[theorem]{Lemma}
\newtheorem{proposition}[theorem]{Proposition}
\newtheorem{corollary}[theorem]{Corollary}
\newtheorem{imported}[theorem]{Imported theorem}
\theoremstyle{definition}
\newtheorem{definition}[theorem]{Definition}
\newtheorem{example}[theorem]{Example}
\theoremstyle{remark}
\newtheorem{remark}[theorem]{Remark}
\newcommand{\N}{\mathbb N}
\newcommand{\Z}{\mathbb Z}
\newcommand{\Q}{\mathbb Q}
\newcommand{\R}{\mathbb R}
\newcommand{\one}{\mathbf 1}
\newcommand{\zero}{\mathbf 0}
\newcommand{\supp}{\operatorname{supp}}
\newcommand{\Max}{\operatorname{Max}}
\newcommand{\Int}{\operatorname{int}}
\newcommand{\conv}{\operatorname{conv}}
\newcommand{\rank}{\operatorname{rank}}
\newcommand{\cov}{\operatorname{cov}}
\newcommand{\Deg}{\operatorname{Deg}}
\newcommand{\Dr}{\mathcal D}
\newcommand{\Ical}{\mathcal I}
\newcommand{\Ecal}{\mathcal E}
\newcommand{\Tr}{\mathcal T}
\newcommand{\Dz}[1]{\Delta^{0}_{#1}}
\newcommand{\qint}[1]{[#1]_q}
\newcommand{\Zeta}{\mathcal Z}
\newcommand{\zt}[1]{\mathbf z^{#1}}
\newcommand{\omitprod}[1]{\prod_{i\notin #1}(1-z_i)}
\newcommand{\qbinom}[2]{\genfrac{[}{]}{0pt}{}{#1}{#2}_{q}}
\newcommand{\gausstbinom}[2]{\genfrac{[}{]}{0pt}{}{#1}{#2}_{t}}
\hypersetup{pdftitle={Support reciprocity and the q-zeta polynomial of preorder lattice-point posets},
 pdfauthor={Merged AI-assisted research draft prepared for Vladimir Reshetnikov},
 pdfsubject={A proposed proof of Athanasiadis--Chapoton Conjecture 4.9, with two independent proofs and evaluations at every negative q-integer}}

\title{\color{Ink}\bfseries Support reciprocity and the \texorpdfstring{$q$}{q}-zeta polynomial\protect\\
of preorder lattice-point posets\\[0.45em]
\large Two independent proofs, every negative \texorpdfstring{$q$}{q}-integer,
and arbitrary height functions}
\author{Merged research draft prepared for Vladimir Reshetnikov\\
\small Mathematical investigation and exposition with AI assistance}
\date{20 September 2026}

\begin{document}
\maketitle
\begin{abstract}
Athanasiadis and Chapoton conjecture a negative-evaluation formula for the
$q$-zeta polynomial of the coordinatewise poset of lattice points in a
preorder polytope. This draft gives a proposed proof of the full
$q$-identity, not only its unweighted specialization, and strengthens it to a
separate identity for every support of a maximal lattice point. The
strengthened statement is proved \emph{twice}, by two arguments that share
nothing beyond two classical imports and the elementary structure of the
poset. The first proof runs through weighted incidence inversion, a
support-restricted bipartite graph, and Ehrhart--Macdonald reciprocity at the
capacities $\one+\one_T$; inside it, the geometric step is available in two
interchangeable forms, one importing Postnikov's bipartite duality and one
deriving the duality it needs from a single application-free Euler identity
for arbitrary bipartite graphs. The second proof uses no incidence matrix, no
matching theorem and no second polytope: it keeps an auxiliary capacity alive
through multivariate Ehrhart reciprocity, kills every non-maximal layer by a
single specialization, and converts the result into a coefficient-transfer
theorem for the involution $\Tr f(z)=(1-z)f(-z/(1-z))$. Only the second proof
reaches every negative $q$-integer $\qint{-r}$, where it produces manifestly
nonnegative integer coefficients with a colored-composition meaning; only the
first reaches arbitrary normalized height functions. We also obtain
multivariate, transversal-matroid and Tutte interpretations, a
characterization of the supports that occur, sharp extreme degrees, explicit
coefficient formulas, worked examples, and three obstructions that bound the
theorem in three different directions. Three independent exact verifiers are
included. The proofs are mathematical, with the classical inputs stated
explicitly; the computations are corroboration rather than a proof
substitute.
\end{abstract}

\noindent\textbf{Status.} This is an AI-assisted research draft containing
complete proposed arguments. It has not been independently refereed or
verified in a proof assistant. The source checked still labels the target as
Conjecture~4.9; focused literature searches found no separate resolution.
This is not a certification of priority. Where the three source drafts merged
here recorded a caveat differently, the most cautious wording has been kept.

\tableofcontents
\clearpage

\section{The problem, the results, and the nonduplication boundary}
\label{sec:problem}

Let $E$ be a finite set of cardinality $n$, and let $\tau$ be a preorder on
$E$: a reflexive, transitive relation, with antisymmetry \emph{not} required.
An \emph{ideal} is a set $I\subseteq E$ such that $j\in I$ and $i\le_\tau j$
imply $i\in I$; a \emph{filter} is the complement of an ideal, equivalently an
upward closed set. For a vector $x$, write $x(I)=\sum_{i\in I}x_i$ and
$|x|=x(E)$. Define
\begin{equation}
 Q_\tau=\{x\in\R_{\ge0}^E:x(I)\le |I|\text{ for every ideal }I\},
 \qquad P_\tau=Q_\tau\cap\Z^E.
 \label{eq:preorder-polytope}
\end{equation}
The order on $P_\tau$ is coordinatewise. Its rank is $|a|$, and all maximal
points have rank $n$; proofs of these basic facts are included below.

Put $\qint m=(q^m-1)/(q-1)$. The $q$-zeta polynomial is characterized by
\begin{equation}
 \Zeta_q(P,\qint m)
 =\sum_{a_1\le\cdots\le a_{m-1}}q^{|a_1|+\cdots+|a_{m-1}|}
 \qquad(m\ge2).
 \label{eq:qzeta-def}
\end{equation}
Here the chains are weak chains, repetitions are allowed, and there are
\emph{$m-1$} entries. The second argument of $\Zeta_q$ is an ordinary
polynomial variable over $\Q(q)$, not a count of chain entries. In particular,
$\qint{-1}=-q^{-1}$. These conventions are those in
Athanasiadis--Chapoton~\cite[Section~4.4]{AC}, following
Chapoton~\cite{Chapoton}.

The target is Conjecture~4.9, equation~(21), in~\cite{AC}:
\begin{equation}
 \boxed{\displaystyle
 \Zeta_q(P_\tau,-q^{-1})
 =(-1)^n\sum_{a\in\Max(P_\tau)}q^{-\cov(a)}.}
 \label{eq:target}
\end{equation}
The statistic $\cov(a)$ counts the elements covered by $a$, not those covering
$a$. For our coordinate downset,
\begin{equation}
 \cov(a)=|\supp(a)|,\qquad \supp(a)=\{i:a_i>0\}.
 \label{eq:covers-support}
\end{equation}
It is this identity \eqref{eq:covers-support}, proved in
Section~\ref{sec:transport}, that makes \eqref{eq:target} literally the source
conjecture rather than a renormalized variant of it.

\subsection{What is proved}
\label{sec:whatisproved}

For independent indeterminates $(z_i)_{i\in E}$, set
\begin{equation}
 H_\tau(\mathbf z)
 =\sum_{a\in\Max(P_\tau)}\prod_{i\in\supp(a)}z_i,
 \qquad \zt B=\prod_{i\in B}z_i,
 \label{eq:support-poly}
\end{equation}
and write $H_\tau(u)=H_\tau(u,\ldots,u)$ for its diagonal. The polynomial
$H_\tau$ is multiaffine: each variable has degree at most one. Different
maximal points may contribute the same monomial. For $B\subseteq E$ put
\begin{align}
 N_\tau(B)&=\#\{a\in\Max(P_\tau):\supp(a)=B\},
 \label{eq:N-def}\\
 \eta_\tau(B)&=\sum_{\substack{S\subseteq E\\\one_B+\one_S\in P_\tau}}
                 (-1)^{|S|}.
 \label{eq:eta-def}
\end{align}

\begin{theorem}[Support-refined reciprocity]
\label{thm:main}
For every finite preorder $\tau$ on $E$,
\begin{equation}
 H_\tau(\mathbf z)
 =(-1)^n
 \sum_{\substack{B,S\subseteq E\\ \one_B+\one_S\in P_\tau}}
 (-1)^{|B|+|S|}\zt B.
 \label{eq:main}
\end{equation}
Equivalently, for every $B\subseteq E$,
\begin{equation}
 N_\tau(B)=(-1)^{n-|B|}\eta_\tau(B)
 =(-1)^{n-|B|}
 \sum_{\substack{S\subseteq E\\\one_B+\one_S\in P_\tau}}(-1)^{|S|}.
 \label{eq:exact-support}
\end{equation}
Conjecture~\eqref{eq:target} follows by setting every $z_i=q^{-1}$ and using
weighted incidence inversion.
\end{theorem}

The second headline statement goes past $\qint{-1}$. For integers $r\ge1$ and
$a\ge0$, define
\begin{equation}
 C_{r,0}(t)=1,\qquad
 C_{r,a}(t)=\sum_{s=1}^{\min(r,a)}
 \binom{a-1}{s-1}e_s(t,t^2,\ldots,t^r)\quad(a\ge1),
 \label{eq:Cdefinition}
\end{equation}
where $e_s$ is the elementary symmetric polynomial of degree $s$. Every
coefficient of $C_{r,a}$ is a nonnegative integer.

\begin{theorem}[All negative $q$-integers]
\label{thm:allnegative}
For every finite preorder $\tau$ and every integer $r\ge1$,
\begin{equation}
 (-1)^n \Zeta_q(P_\tau,\qint{-r})
 =\sum_{b\in \Max(P_\tau)}\prod_{i\in E}C_{r,b_i}(q^{-1})
 \ \in\ \Z_{\ge0}[q^{-1}].
 \label{eq:allnegative}
\end{equation}
For $n>0$, the polynomial on the right, in the variable $t=q^{-1}$, is monic
of degree $rn$; its smallest exponent is the number $c_\tau$ of maximal
equivalence classes of $\tau$, and the coefficient there is the number of
maximal points of that support size.
\end{theorem}

Here an equivalence class is a class of the relation $i\sim j$ if $i\le_\tau
j\le_\tau i$; the quotient is a finite poset. Equivalence classes retain their
original cardinalities throughout: replacing a class by a single element and
forgetting that cardinality would change the problem. At $r=1$ we have
$C_{1,a}(t)=t$ for $a\ge1$, so Theorem~\ref{thm:allnegative} contains
\eqref{eq:target}.

The third headline statement goes sideways instead, replacing the rank by an
arbitrary height statistic. Let $h:P_\tau\to\N$ be strictly increasing on
comparable distinct elements with $h(\zero)=0$, and let $\Zeta_{q,h}$ be the
height-weighted $q$-zeta polynomial of~\cite{Chapoton}, defined by
\eqref{eq:qzeta-def} with $|a_j|$ replaced by $h(a_j)$.

\begin{theorem}[Arbitrary normalized heights]
\label{thm:height-intro}
For every such $h$,
\begin{equation}
 \Zeta_{q,h}(P_\tau,\qint{-1})
 =(-1)^n\sum_{b\in\Max(P_\tau)}q^{-h(\one_{\supp(b)})}.
 \label{eq:height-intro}
\end{equation}
The exponent uses the \emph{Boolean support vector} $\one_{\supp(b)}$, not
$h(b)$.
\end{theorem}

\subsection{Two proofs, and two engines inside the first}
\label{sec:roadmap}

Theorem~\ref{thm:main} is proved twice. The two proofs are recorded in full,
and deliberately so: the result is unrefereed, the arguments are independent,
and each reaches a generalization the other does not.

\medskip\noindent\textbf{Proof A (incidence inversion plus geometry),
Sections~\ref{sec:incidence}--\ref{sec:private}.} Weighted incidence-matrix
inversion over $\Q(q)$ reduces the negative evaluation to points with
coordinates in $\{0,1,2\}$ and produces the signed Boolean sums
$\eta_\tau(B)$; this step works for an arbitrary finite coordinate downset and
for an arbitrary normalized height, which is where
Theorem~\ref{thm:height-intro} comes from. A geometric argument then evaluates
those signed sums. Two interchangeable engines are available for it:
\begin{itemize}
\item \textbf{Engine~1 (Section~\ref{sec:interior}), cumulative supports.}
 A support-restricted bipartite graph $G_T$ and Postnikov's bipartite
 draconian duality identify maximal points supported inside $T$ with the
 interior lattice points of $Q_{\tau^*}(\one+\one_T)$; Ehrhart--Macdonald
 reciprocity and a coordinatewise two-value interpolation finish. Shortest
 route, largest imported surface.
\item \textbf{Engine~2 (Section~\ref{sec:private}), exact supports.}
 Two Hall lemmas reduce $\eta_\tau(B)$ at the \emph{exact} support $B$ to an
 alternating count over the independent sets of a private-element bipartite
 graph, and a standalone Euler identity, proved for \emph{every} bipartite
 graph from two applications of Postnikov's lattice-count theorem plus
 Ehrhart, evaluates that count. Longest route, smallest imported surface: it
 re-derives, in the case it needs, the duality Engine~1 imports.
\end{itemize}

\medskip\noindent\textbf{Proof B (the transfer involution),
Section~\ref{sec:transfer}.} No incidence matrix, no Hall theorem, no
bipartite graph, no second polytope, no interior-point count. Instead: build
the independent-capacity enumerator $\Ecal_\sigma(u,v)$ and prove
$\Ecal_\sigma(-u,-v)=(-1)^n\Ecal_\sigma(u-\one,v-1)$ by interior translation
and Ehrhart; then set $v=0$, so that the auxiliary capacity slides from $0$ to
$-1$ and the binomial factor $\binom{n-|a|-1}{n-|a|}$ annihilates every
non-maximal layer; then convert the resulting binomial identity into a
transfer theorem for arbitrary local weights, through the involution
$\Tr f(z)=(1-z)f(-z/(1-z))$ on the basis $(1-z)^m$; finally evaluate
$\Zeta_q$ by an explicit Gaussian endpoint formula valid at \emph{every}
integer $m$ and feed $f(z)=(1-z)(1-q^{-1}z)$ through $\Tr$. This is the only
route that reaches $\qint{-r}$ for every $r$, so it is where
Theorem~\ref{thm:allnegative} comes from.

\medskip
The two proofs share the two classical imports (Postnikov's lattice-count
theorem and Ehrhart--Macdonald reciprocity), the elementary structure of
$P_\tau$ in Section~\ref{sec:transport}, and nothing else. They also touch at
two seams, which Sections~\ref{sec:transfer} and~\ref{sec:consequences} point
out explicitly: the two-capacity collapse of Proof~B is the two-value
interpolation table of Proof~A in binomial rather than polynomial language,
and the independent-set family of Engine~2 is the family
$\{S:\one_B+\one_S\in P_\tau\}$ of Engine~1, which is a transversal matroid.

\subsection{What is deliberately kept more than once}
\label{sec:kept}

Several things below exist in two or three versions. None of the duplications
is an accident of merging: each buys something the others do not, and the
place where each lives is listed here once and for all.

\begin{enumerate}[label=(\arabic*),leftmargin=2.2em]
\item \textbf{Two proofs of Theorem~\ref{thm:main}}
 (Sections~\ref{sec:incidence}--\ref{sec:private} and
 Section~\ref{sec:transfer}): they cross-validate an unrefereed result, and
 only one of them generalizes to $\qint{-r}$ while only the other reaches
 arbitrary height functions.
\item \textbf{Two geometric engines inside Proof~A}
 (Sections~\ref{sec:interior} and~\ref{sec:private}): Engine~1 is much
 shorter, Engine~2 imports strictly less, so a reader who declines
 Imported theorem~\ref{imp:duality} still has a complete route.
\item \textbf{Three derivations of the exact-support identity}
 \eqref{eq:exact-support} (Sections~\ref{sec:completion}, \ref{sec:private}
 and~\ref{sec:proofB-support}): this is the statement the whole document is
 built on, so every independent route to it is worth having.
\item \textbf{Two constructions of the $q$-zeta polynomial at negative
 arguments} (Lemma~\ref{lem:continuation} and
 Proposition~\ref{prop:endpoint}): the matrix continuation is what makes
 arbitrary heights reachable, the endpoint formula is what makes every
 $\qint{-r}$ reachable, and neither implies the other
 (Remark~\ref{rem:two-constructions}).
\item \textbf{Three obstructions carrying four failure statements}
 (Section~\ref{sec:limits}): they bound purity, transitivity and unit
 capacities separately, and collapsing any of them would lose a boundary.
\item \textbf{Two matroid invariants for the same coefficient}
 (Proposition~\ref{prop:matroid} and \eqref{eq:h-polynomial}): the Tutte
 evaluation places it inside a standard formalism, the independence-complex
 $h$-polynomial needs no matroid theorem at all.
\item \textbf{Three verifiers} (Section~\ref{sec:computation}): they test
 three different theorems, at two different seeds and two different sampling
 depths, so the sampled populations are genuinely independent.
\item \textbf{Two statements about which supports occur}
 (Proposition~\ref{prop:support-feasible} and
 Lemma~\ref{lem:minsupp} with Proposition~\ref{prop:degrees}): the
 qualitative cofinality criterion settles positivity, the quantitative
 extremal statement is what the degree claim of
 Theorem~\ref{thm:allnegative} depends on.
\item \textbf{Two sets of coefficient formulas}
 (Proposition~\ref{prop:top} and Section~\ref{sec:consequences}): the
 third-from-the-top coefficient \eqref{eq:next-coefficient} and the product
 rule over disjoint unions came from different arguments and neither follows
 from the other.
\item \textbf{Two accounts of the $q\to1$ specialization}
 (Section~\ref{sec:shadow}): the nilpotent argument is the cleaner
 justification, the interpolation is what the verifier implements, and both
 are needed to explain why the case already covered by
 \cite{PriorReciprocity} falls out safely rather than being reproved.
\item \textbf{Three reciprocity references}
 (Section~\ref{sec:imports-ehrhart}): historical, textbook and
 cone-reciprocity, for readers arriving from three different literatures.
\item \textbf{Two seams between the two proofs}
 (Remarks~\ref{rem:seam-table} and~\ref{rem:seam-matroid}): the two-capacity
 collapse against the two-value table, and Engine~2's independent sets
 against Engine~1's transversal matroid; recording them is what keeps the
 coincidences from looking like repetitions.
\end{enumerate}

\subsection{Provenance}
\label{sec:provenance}

This document is a merge of three separate research drafts prepared on
20~September 2026, each of which independently proposed a proof of
\eqref{eq:target} \emph{and}, letter for letter, of the exact-support
strengthening \eqref{eq:exact-support}. All three are directories of the same
collection:
\begin{description}[style=nextline,leftmargin=1.2em,labelwidth=0pt,labelindent=0pt,font=\normalfont\bfseries]
\item[\nolinkurl{enumerative-combinatorics/preorder-q-zeta-support-reciprocity}]
 The spine of the present document. Proved \eqref{eq:main} in its multivariate
 form by weighted incidence inversion, the support--interior identity at
 capacities $\one+\one_T$ through Postnikov's bipartite duality, univariate
 Ehrhart reciprocity along each capacity ray, a two-value interpolation table
 and a Bernstein-type collapse. Contributed Engine~1, the Tutte evaluation,
 the product rule, the nontransitive and capacity-two obstructions, and most
 of the package infrastructure.
\item[\nolinkurl{enumerative-combinatorics/preorder-q-zeta-supportwise}]
 Proved \eqref{eq:exact-support} at the exact support by residual and
 privatization Hall lemmas, followed by a self-contained Euler identity for
 arbitrary bipartite graphs obtained from two applications of Postnikov's
 Theorem~11.3. Contributed Engine~2, the arbitrary-height theorem, the
 $[t^{n-2}]$ coefficient formula, the finite Hall proof, the standalone
 bipartite test suite, and the pure-downset obstruction in its supportwise
 form.
\item[\nolinkurl{enumerative-combinatorics/preorder-q-reciprocity-weighted}]
 Proved \eqref{eq:target} and \eqref{eq:exact-support} by independent-capacity
 reciprocity, the $v=0$ pivot, and the coefficient-transfer involution, with
 no geometric duality and no matching theorem at all. Contributed Proof~B,
 Theorem~\ref{thm:allnegative} and the whole $\qint{-r}$ apparatus, the
 Gaussian endpoint formula, the colored-composition model, the sharp extreme
 degrees, and the pure-downset obstruction in its transfer form.
\end{description}
The merge states the shared theorem once, keeps both proofs, keeps both
engines inside Proof~A, and keeps every statement that exists in only one of
the three. Section~\ref{sec:audit} records which caveat came from where, and
Appendix~\ref{app:ledger} merges the three source ledgers.

\subsection{Relation to the other packages in this collection}
\label{sec:nonduplication}

The supplied manifest~\cite{Manifest} includes a package entitled
\emph{Reciprocity and Matrix Duality for Preorder Polytopes}, in the directory
\nolinkurl{enumerative-combinatorics/preorder-polytope-reciprocity}; we cite
it as~\cite{PriorReciprocity}. Its description explicitly claims the ordinary
$q=1$ case of Conjecture~4.9 and explicitly excludes the full
$q$-refinement. The present target is exactly that excluded refinement. The
manifest also contains results about other conjectures of the same source, and
a separate entry, \emph{A Reflexive Root-Polytope Model for Preorder
$h$-Polynomials}, which concerns a different conjecture and is not repeated
here.

The difference is substantive. At $q=1$, the right side of \eqref{eq:target}
retains only the total number of maximal points. At general $q$, it records
the entire distribution of their support sizes. Theorem~\ref{thm:main} retains
the actual support set, which is finer still, and
Theorem~\ref{thm:allnegative} leaves the point $\qint{-1}$ altogether.

Three specific overlaps with \cite{PriorReciprocity} are acknowledged rather
than concealed, and are flagged again at the places where they occur.
\begin{enumerate}[label=(\roman*)]
\item The independent-capacity reciprocity of
 Proposition~\ref{prop:reciprocity} is the identity already recorded in the
 manifest entry for~\cite{PriorReciprocity}, whose Theorem on multivariate
 reciprocity states it. We reprove it here in order to make Proof~B
 independent of an unrefereed draft, not because it is claimed as new.
\item The enumerator formula \eqref{eq:L-formula} for $L_{\tau^*}$ is the
 specialization at $s=0$ of the rising-factorial enumerator of
 \cite{PriorReciprocity}, derived there the same way from Postnikov's theorem
 and the filter identification of draconian vectors, with both drafts
 crediting \cite[Lemma~4.1]{AC} for that identification. It is re-derived here
 and is not claimed as a new theorem.
\item The ordinary case $\Zeta(P_\tau,-1)=(-1)^n\#\Max(P_\tau)$, recovered in
 Section~\ref{sec:shadow}, is the corollary already proved in
 \cite{PriorReciprocity}. We continue to present it as covered there and not
 as a fresh result; what is new here is the $q$-refinement above it and the
 support refinement beyond that.
\end{enumerate}
No proof in \cite{PriorReciprocity} is imported as a lemma, and no uninspected
proof in that package is assumed correct. Conversely, nothing here supersedes
it. Only the manifest, and not the full text of the earlier packages, was
available to the three drafts when they assessed nonduplication; that
limitation is retained as stated.

Finally, a nonduplication boundary inside the source paper itself:
$H_\tau$ counts \emph{maximal} points by support, whereas the preorder
$h$-polynomial of \cite[Section~5]{AC} counts support sizes over \emph{all}
lattice points. These are different polynomials, and no claim is made about
the latter.

There are three explicitly identified classical ingredients: the integral
bipartite matching criterion (Hall's theorem, proved in
Appendix~\ref{app:hall}), Ehrhart--Macdonald reciprocity, and Postnikov's
lattice-count and bipartite-duality theorems. The steps relating these
ingredients to support sets, and their combination with the negative
$q$-evaluation, are proved in full. No conjecture about real-rootedness,
unimodality, or the $h$-polynomials of preorder polytopes is used.

\section{Preorder polytopes as transport polytopes}
\label{sec:transport}

For $J\subseteq E$, let
\[
 D_\tau(J)=\{i:i\le_\tau j\text{ for some }j\in J\},\qquad
 U_\tau(J)=\{i:j\le_\tau i\text{ for some }j\in J\}.
\]
Transitivity makes these an ideal and a filter, respectively. Reflexivity
gives $J\subseteq D_\tau(J)$ and $J\subseteq U_\tau(J)$. We abbreviate the
closure of a singleton in the usual way. The dual preorder $\tau^*$ reverses
all comparisons, so that the ideals of $\tau^*$ are the filters of $\tau$.

For a nonnegative capacity vector $\mathbf t=(t_i)$ and a nonnegative scalar
$v$, define
\begin{equation}
\begin{aligned}
 Q_\tau(\mathbf t)&=
 \{x\ge0:x(I)\le t(I)\ \text{ for every ideal }I\},\\
 K_\tau(\mathbf t,v)&=
 \{x\ge0:x(I)\le t(I)+v\ \text{ for every nonempty ideal }I\}.
\end{aligned}
 \label{eq:weighted-Q}
\end{equation}
Thus $Q_\tau=Q_\tau(\one)$ and $Q_\tau(\mathbf t)=K_\tau(\mathbf t,0)$. The
auxiliary capacity $v$ plays no role in Proof~A; it is the whole mechanism of
Proof~B, and it is introduced here once so that both proofs can use the same
Minkowski realization.

\begin{lemma}[Minkowski realization and transport description]
\label{lem:transport}
Write $\Dz A=\conv(\{\zero\}\cup\{e_j:j\in A\})$, so that
$\Dz\emptyset=\{\zero\}$. For all nonnegative real capacities,
\begin{equation}
 K_\tau(\mathbf t,v)
 =\sum_{i\in E} t_i\,\Dz{U_\tau(i)}\ +\ v\,\Dz E .
 \label{eq:minkowski}
\end{equation}
For integral $\mathbf t$ and $v$, this is a lattice polytope. If $\mathbf t>0$
coordinatewise, then $Q_\tau(\mathbf t)=K_\tau(\mathbf t,0)$ has dimension
$n$; so does $K_\tau(\mathbf t,v)$ whenever $v>0$.

Moreover $Q_\tau(\mathbf t)$ is the transport polytope in which each $i\in E$
is a source supplying at most $t_i$ units, each $j\in E$ is a destination
requiring $x_j$ units, transport from $i$ to $j$ is allowed exactly when
$i\le_\tau j$, and unused supply is permitted; equivalently,
\begin{equation}
 x\in Q_\tau(\mathbf t)
 \iff
 x\ge0\ \text{and}\
 x(J)\le t(D_\tau(J))\ \ \text{for every }J\subseteq E.
 \label{eq:hall-demand}
\end{equation}
\end{lemma}

\begin{proof}
Denote the Minkowski sum in \eqref{eq:minkowski} by $M$. A summand
$t_i\Dz{U_\tau(i)}$ contributes at most $t_i$ to an ideal $I$, and contributes
zero if $i\notin I$, since an element of $U_\tau(i)\cap I$ would force
$i\in I$. The last summand contributes at most $v$. Hence
$M\subseteq K_\tau(\mathbf t,v)$.

For the reverse inclusion, fix $x\in K_\tau(\mathbf t,v)$ and a linear
functional $w\in\R^E$. Put $w^+_j=\max(w_j,0)$ and $J_s=\{j:w_j\ge s\}$ for
$s>0$. When $J_s$ is nonempty, its downward closure $D_\tau(J_s)$ is a
nonempty ideal, so $x(J_s)\le x(D_\tau(J_s))\le t(D_\tau(J_s))+v$.
Integrating this step-function inequality,
\begin{align*}
 w\cdot x\le w^+\cdot x=\int_0^\infty x(J_s)\,ds
 \le\sum_{i\in E}t_i\max\bigl(0,\max_{j\in U_\tau(i)}w_j\bigr)
       +v\max\bigl(0,\max_{j\in E}w_j\bigr),
\end{align*}
and the right-hand side is exactly the support function of $M$ at $w$. Both
sets are compact, because nonnegativity and the constraint for $I=E$ bound
all coordinates; the separation theorem gives $x\in M$. A finite Minkowski
sum of integer multiples of lattice simplices is a lattice polytope.
Reflexivity puts the box $\prod_i[0,t_i]$ inside $Q_\tau(\mathbf t)$, which
has nonempty interior when every $t_i>0$; and $v\Dz E$ is full-dimensional
for $v>0$.

For the transport description, any transport satisfies the ideal
inequalities, because a unit received by a destination in an ideal $I$ must
have originated in $I$. Conversely, the capacitated bipartite matching
criterion says that the demands can be met if and only if
\eqref{eq:hall-demand} holds: these are the source-capacity bounds for the
neighbours of each set of destinations, that is, the cut conditions of the
corresponding flow network. The ideal inequalities imply them, since
$x(J)\le x(D_\tau(J))\le t(D_\tau(J))$, and conversely every ideal $I$
satisfies $D_\tau(I)=I$. The matching criterion itself is Hall's theorem,
proved in Appendix~\ref{app:hall}; the passage from vector inequalities to an
ordinary finite matching replaces a demand of multiplicity $x_j$ by $x_j$
distinct copies, as explained there.
\end{proof}

We restate the unit-capacity case in the language of allocations, because
Proof~A's Engine~2 uses it constantly.

\begin{lemma}[Allocation criterion]
\label{lem:allocation}
Think of each $r\in E$ as one available resource and of $a_i$ as a number of
demands of type $i$, a resource $r$ being able to serve a demand of type $i$
exactly when $r\le_\tau i$. For $a\in\N^E$ the following are equivalent:
\emph{(i)} $a\in P_\tau$; \emph{(ii)} the $|a|$ demands can be assigned
injectively to resources along allowed edges; \emph{(iii)}
$a(A)\le|D_\tau(A)|$ for every $A\subseteq E$.
\end{lemma}

\begin{proof}
This is \eqref{eq:hall-demand} at $\mathbf t=\one$, together with the
integrality of the matching: (i)~$\Rightarrow$~(iii) since
$a(A)\le a(D_\tau(A))\le|D_\tau(A)|$; (iii)~$\Rightarrow$~(ii) by Hall's
theorem applied to $a_i$ distinct copies of each demand $i$, whose possible
resources form $D_\tau(A)$ where $A$ is the set of their types; and
(ii)~$\Rightarrow$~(i) because every demand whose type lies in an ideal $I$
must use a resource in $I$.
\end{proof}

\begin{lemma}[Maximal points and intervals]
\label{lem:max}
The set $P_\tau$ is a finite coordinate downset containing $\zero$ and
$\one$, indeed containing every $\one_S$. Its maximal elements are exactly its
points of total coordinate sum $n$. Every interval $[a,b]$ is the coordinate
box $\prod_i[a_i,b_i]\cap\Z^E$. In particular $P_\tau$ is graded with rank
$|a|$, is pure of rank $n$, and \eqref{eq:covers-support} holds.
\end{lemma}

\begin{proof}
Downward closure and the inclusion of the Boolean cube follow directly from
the inequalities, since $|S\cap I|\le|I|$. The inequality for $I=E$ gives
$|a|\le n$, proving finiteness. If $a$ is integral with $|a|<n$, choose an
integral allocation as in Lemma~\ref{lem:allocation}; at least one resource
$r$ is unused, and adding one demand of type $r$ served by $r$ itself, which
reflexivity permits, gives $a+e_r\in P_\tau$. Thus $a$ is not maximal, while a
point of total $n$ cannot be increased in any coordinate.

If $a\le c\le b$ and $b\in P_\tau$, downward closure gives $c\in P_\tau$,
which proves the interval assertion. Hence the elements covered by $a$ are
exactly the $a-e_i$ with $a_i>0$, which is \eqref{eq:covers-support}, and
decrementing to $\zero$ proves the rank formula.
\end{proof}

\begin{remark}[An alternative proof of the maximal layer]
\label{rem:tight}
The maximal-layer statement can also be proved without any matching theorem,
which matters for Proof~B: if $a\in P_\tau$ has $|a|<n$ and $a+e_i\notin
P_\tau$ for every $i$, then for each $i$ some ideal $I_i\ni i$ is
\emph{tight}, that is $a(I_i)=|I_i|$. Modularity of $x\mapsto x(I)$ and of
cardinality gives
$a(I\cup J)+a(I\cap J)=|I\cup J|+|I\cap J|$ for tight $I,J$, and since each
term on the left is at most the corresponding cardinality, both $I\cup J$ and
$I\cap J$ are tight. The union of the $I_i$ is $E$, so $E$ is tight,
contradicting $|a|<n$.
\end{remark}

\begin{remark}[Preorders are genuinely allowed]
An equivalence class of $\tau$ is not contracted to a single coordinate. Its
individual elements remain distinct unit-capacity sources and distinct
destinations. Ideals contain a whole equivalence class or none of it, but
support sets need not. This distinction is important both in the proofs and
in the exhaustive tests. For the empty ground set,
$P_\tau=\Max(P_\tau)=\{()\}$ and all the polynomials below equal $1$; the
geometric statements are proved for $n\ge1$ and this convention is used at
$n=0$.
\end{remark}

\section{Negative evaluation by weighted incidence inversion}
\label{sec:incidence}

This section is the first half of Proof~A. Everything in it is valid for an
arbitrary finite poset with least element and an arbitrary normalized height,
hence in particular for any finite coordinate downset $P\subseteq\N^E$, not
just a preorder polytope. It is a specialization of the $q$-incidence
viewpoint of Chapoton~\cite[Appendix~C]{Chapoton}, with the matrix conventions
and the negative exponent verified explicitly.

\subsection{An algebraic continuation, not an analytic one}

Let $P$ be a finite poset with least element $\zero$, and let $h:P\to\N$ be
\emph{normalized}: strictly increasing on comparable distinct elements, with
$h(\zero)=0$. Order the elements by nondecreasing $h$. Let $Z$ be the
incidence zeta matrix, let $D$ be the weight matrix, and put
\begin{equation}
 Z_{a,b}=\begin{cases}1&a\le b,\\0&\text{otherwise},\end{cases}
 \qquad D_{a,a}=q^{h(a)},\qquad A=ZD.
 \label{eq:matrices}
\end{equation}
Use row vectors on the left: write $e_0^{\mathsf T}$ for the row supported at
$\zero$ and $\mathbf j$ for the column vector of ones; $\mathbf j$ is not an
element of $P$. For every integer $\ell\ge0$,
\begin{equation}
 e_0^{\mathsf T}A^{\ell}\mathbf j
 =\sum_{a_1\le\cdots\le a_{\ell}}
   q^{h(a_1)+\cdots+h(a_{\ell})},
 \label{eq:positive-matrix}
\end{equation}
the empty-chain sum at $\ell=0$ being $1$. Comparing with
\eqref{eq:qzeta-def} at $\ell=m-1$ identifies \eqref{eq:positive-matrix} with
$\Zeta_{q,h}(P,\qint m)$ for $m\ge2$.

\begin{lemma}[Polynomial continuation to negative integers]
\label{lem:continuation}
Let $g_0<\cdots<g_d$ be the distinct values of $h$. The expression
$e_0^{\mathsf T}A^{m-1}\mathbf j$, for all integers $m$, is a polynomial of
degree at most $g_d$ in $\qint m$, with coefficients in $\Q(q)$. Consequently
\begin{equation}
 \Zeta_{q,h}(P,-q^{-1})=e_0^{\mathsf T}A^{-2}\mathbf j
 =e_0^{\mathsf T}\mu D^{-1}\mu\,\mathbf j,
 \qquad \mu=Z^{-1}.
 \label{eq:negative-matrix}
\end{equation}
\end{lemma}

\begin{proof}
Let $V_j$ be the space of row vectors supported on elements of height at least
$g_j$. Right multiplication by $A$ preserves $V_j$, and on $V_j/V_{j+1}$ it
acts as the scalar $q^{g_j}$, because two comparable elements of the same
height are equal. Therefore $A-q^{g_j}I$ maps $V_j$ into $V_{j+1}$, and
applying the factors in increasing order gives
\[
 \prod_{j=0}^{d}(A-q^{g_j}I)=0 .
\]
Over $\Q(q)$ the factors are distinct, so $A$ is diagonalizable with nonzero
eigenvalues among $q^{g_0},\ldots,q^{g_d}$. Hence there are $c_j\in\Q(q)$
with
\[
 e_0^{\mathsf T}A^{\ell}\mathbf j=\sum_{j=0}^{d}c_jq^{g_j\ell}
 \qquad\text{for \emph{every} integer }\ell .
\]
Put $H(X)=\sum_jc_jq^{-g_j}X^{g_j}$. For $m\ge2$ and $\ell=m-1$ this equals
$H(q^m)=H(1+(q-1)\qint m)$, and such a polynomial is unique because the
$\qint m$, $m\ge2$, are infinitely many distinct elements of $\Q(q)$. So
$\Zeta_{q,h}(P,t)=H(1+(q-1)t)$, and at $t=\qint{-1}=-q^{-1}$ the argument of
$H$ is $q^{-1}$, that is $\ell=-2$: at $m=-1$ the matrix exponent is $-2$,
not $-1$, precisely because the multichains in \eqref{eq:qzeta-def} have
$m-1$ entries. Finally $A^{-1}=D^{-1}\mu$ and $e_0^{\mathsf T}D^{-1}=e_0^{\mathsf T}$
because $h(\zero)=0$. No limiting or analytic-continuation assertion is
involved anywhere in this argument; every step is an identity of polynomials
over $\Q(q)$.
\end{proof}

\begin{remark}
The argument just given also supplies, in this least-element setting, the
polynomial existence that Chapoton's construction \cite{Chapoton} requires;
we do not assume it.
\end{remark}

\subsection{Boolean localization}

From now on $P\subseteq\N^E$ is a finite coordinate downset and $h$ is
arbitrary normalized as above. Its intervals are products of integer chains by
Lemma~\ref{lem:max}. A chain interval of positive length has M\"obius number
$-1$ at length one and $0$ at greater length, directly from the defining
recurrence $\sum_{a\le c\le b}\mu(a,c)=\delta_{a,b}$. Multiplying over
coordinates,
\begin{equation}
 \mu(a,b)=
 \begin{cases}
 (-1)^{|T|},&b-a=\one_T\text{ for some }T\subseteq E,\\
 0,&\text{otherwise}.
 \end{cases}
 \label{eq:mobius}
\end{equation}
In particular $\mu(\zero,a)$ vanishes unless $a$ is Boolean. This
localization is why a statistic of full maximal vectors can be recovered from
vectors with coordinates only $0,1,2$.

\begin{proposition}[The Boolean localization formula]
\label{prop:local}
For every nonempty finite coordinate downset $P\subseteq\N^E$ and every
normalized height $h$,
\begin{equation}
 \Zeta_{q,h}(P,-q^{-1})
 =\sum_{B\subseteq E}(-1)^{|B|}q^{-h(\one_B)}\eta_P(B),
 \qquad
 \eta_P(B)=\sum_{\substack{S\subseteq E\\\one_B+\one_S\in P}}(-1)^{|S|}.
 \label{eq:reduction}
\end{equation}
For the ordinary rank $h=|\cdot|$ this reads
\begin{equation}
 \Zeta_q(P,-q^{-1})
 =\sum_{\substack{B,S\subseteq E\\\one_B+\one_S\in P}}
   (-1)^{|B|+|S|}q^{-|B|},
 \label{eq:local}
\end{equation}
and, with $u=q^{-1}$, equivalently
\begin{equation}
 \Zeta_q(P,-q^{-1})
 =\sum_{c\in P\cap\{0,1,2\}^E}
   (-1)^{s(c)}(1+u)^{s(c)}u^{d(c)},
 \label{eq:ternary}
\end{equation}
where $s(c)=\#\{i:c_i=1\}$ and $d(c)=\#\{i:c_i=2\}$.
\end{proposition}

\begin{proof}
Expand \eqref{eq:negative-matrix}:
\[
 e_0^{\mathsf T}\mu D^{-1}\mu\,\mathbf j
 =\sum_{b\le c}\mu(\zero,b)q^{-h(b)}\mu(b,c)
 =\sum_{a\in P}\mu(\zero,a)q^{-h(a)}\sum_{c\ge a}\mu(a,c).
\]
By \eqref{eq:mobius} a surviving $b$ is $\one_B$ and a surviving $c-b$ is
$\one_S$; conversely, if $\one_B+\one_S\in P$ then downward closure puts both
points in the sum. This is \eqref{eq:reduction}, and \eqref{eq:local} is its
specialization. For \eqref{eq:ternary}, group the pairs $(B,S)$ by the vector
$c=\one_B+\one_S$: a coordinate of $c$ equal to zero contributes $1$; a
coordinate equal to one arises in two ways and contributes $-1-u$; a
coordinate equal to two contributes $u$. Multiply over coordinates.
\end{proof}

Thus \eqref{eq:main} implies \eqref{eq:target}, since the two signs multiply
to $(-1)^n$ and \eqref{eq:covers-support} identifies support size with the
number of lower covers; and \eqref{eq:main} together with \eqref{eq:reduction}
implies Theorem~\ref{thm:height-intro}. What remains for Proof~A is the
evaluation of the signed Boolean sums $\eta_\tau(B)=\eta_{P_\tau}(B)$. Neither
the matrix calculation nor the Boolean M\"obius formula alone proves the
desired sign or positivity; that is the work of
Sections~\ref{sec:imports}--\ref{sec:private}.

The rational function in \eqref{eq:local} is a Laurent polynomial. Although
the full poset may contain large coordinate values, this particular evaluation
sees only its ternary-coordinate portion. This fact alone does \emph{not}
identify the answer with maximal points: Section~\ref{sec:limits} exhibits
downsets where it is false. The special preorder geometry is still needed.

\section{The classical polyhedral inputs, with conventions fixed}
\label{sec:imports}

We state the precise versions used, so that no normalization of a trimmed
polytope or of a draconian sequence remains implicit. These results are
classical inputs, not contributions of this draft. Both proofs use
Section~\ref{sec:imports-lattice} and Section~\ref{sec:imports-ehrhart};
only Engine~1 of Proof~A uses Section~\ref{sec:imports-duality}.

Throughout, for an indeterminate $X$ and an integer $a\ge0$,
\begin{equation}
 \binom Xa=\frac{X(X-1)\cdots(X-a+1)}{a!},\qquad\binom X0=1,
\end{equation}
\begin{equation}
 f_a(X)=\binom{X+a-1}{a}=\frac{X(X+1)\cdots(X+a-1)}{a!},
 \qquad f_a(X+1)=\binom{X+a}{a}.
 \label{eq:binomial-polynomial}
\end{equation}
Every later evaluation with a negative upper argument means a polynomial
evaluation, not the cardinality of a negatively sized set.

\subsection{Hall's theorem}
The finite marriage theorem \cite{Hall} is used in
Lemmas~\ref{lem:transport}, \ref{lem:allocation}, \ref{lem:euler},
\ref{lem:residual}, \ref{lem:private} and Proposition~\ref{prop:matroid}. It
is stated and proved, with the exact matching convention and the multiplicity
reduction fixed, in Appendix~\ref{app:hall}. It is entirely absent from
Proof~B.

\subsection{The lattice-count formula}
\label{sec:imports-lattice}

\begin{imported}[Postnikov's untrimmed lattice-count formula]
\label{imp:lattice}
Let $F$ be a finite set with $|F|=d+1$, let $I_0=F$, and let
$I_1,\ldots,I_m$ be nonempty subsets of $F$. Write
$\Delta_I=\conv\{e_j:j\in I\}$. For nonnegative integers $y_0,\ldots,y_m$,
\begin{equation}
 \#\left(\left(\sum_{j=0}^{m}y_j\Delta_{I_j}\right)\cap\Z^{F}\right)
 =\sum_{a}\ f_{a_0}(y_0+1)\prod_{j=1}^{m}f_{a_j}(y_j),
 \label{eq:postnikov}
\end{equation}
where $a\in\N^{m+1}$ ranges over the vectors with
\begin{equation}
 \sum_{j=0}^{m}a_j=d,\qquad
 \sum_{j\in J}a_j\le\Bigl|\bigcup_{j\in J}I_j\Bigr|-1
 \quad(\emptyset\ne J\subseteq\{0,\ldots,m\}).
 \label{eq:draconian-conditions}
\end{equation}
\end{imported}

This is \cite[Theorem~11.3]{Postnikov} in its stated untrimmed form; compare
\cite[Theorem~2.1]{AC}. In Postnikov's trimmed formulation, passing to the
untrimmed polytope replaces the coefficient of the full simplex by one more
than that coefficient, which is why the distinguished factor is
$f_{a_0}(y_0+1)=\binom{y_0+a_0}{a_0}$ and not $f_{a_0}(y_0)$. The theorem is
used at zero weights as well, as its nonnegative-integer statement permits:
an auxiliary full simplex may be introduced with $y_0=0$, and its factor is
then $f_{a_0}(1)=1$ for every $a_0\ge0$, \emph{not} $\binom{a_0-1}{a_0}$.
Forgetting this one-unit shift would invalidate every application below.
The imported statement is Postnikov's Theorem~11.3 together with its main
proof; the passage of \cite{Postnikov} explicitly called an ``alternative
semiproof'' is not what is imported here.

\subsection{Bipartite draconian sequences and their duality}
\label{sec:imports-duality}

Let $G=(L,R;\mathcal E)$ be a connected bipartite graph with both parts
nonempty. For $J\subseteq L$ write $N_G(J)\subseteq R$ for its neighbours, and
define
\begin{equation}
 \Dr_L(G)=\left\{a\in\N^L:
 \sum_{i\in L}a_i=|R|-1,\quad
 \sum_{i\in J}a_i\le |N_G(J)|-1\ (\emptyset\ne J\subseteq L)\right\}.
 \label{eq:draconian}
\end{equation}
The set $\Dr_R(G)$ is defined by exchanging the parts.

\begin{imported}[Postnikov's bipartite duality]
\label{imp:duality}
For such a graph, $|\Dr_L(G)|=|\Dr_R(G)|$.
\end{imported}

This is the draconian-sequence form of \cite[Lemma~11.7 and
Corollary~11.8]{Postnikov}. One may view both sets as the lattice points of
two dual trimmed generalized permutohedra; the triangulation interpretation
explains why their cardinalities agree even when their ambient dimensions
differ. Our application in Section~\ref{sec:interior} checks connectedness; it
does not silently assume the theorem for a graph with isolated vertices.

\begin{remark}[The import surface is a real trade]
\label{rem:import-trade}
Imported theorem~\ref{imp:duality} is used only by Engine~1. Engine~2, in
Section~\ref{sec:private}, proves the instance of this duality that it needs
from Imported theorem~\ref{imp:lattice} alone, applied twice to two polytopes
in different ambient dimensions. Engine~1 is much shorter; Engine~2 rests on a
strictly smaller imported surface. A reader who is unwilling to assume
Lemma~11.7 of \cite{Postnikov} still has a complete route to
Theorem~\ref{thm:main}, and this is why both engines are written out.
\end{remark}

\subsection{Ehrhart--Macdonald reciprocity}
\label{sec:imports-ehrhart}

For a full-dimensional lattice polytope $K\subseteq\R^d$, let
$L_K(s)=\#(sK\cap\Z^d)$ for nonnegative integers $s$. Its polynomial
continuation satisfies
\begin{equation}
 L_K(-1)=(-1)^{d}\#(\Int(K)\cap\Z^{d}).
 \label{eq:ehrhart-reciprocity}
\end{equation}
We use this standard reciprocity theorem in precisely this univariate form.
The historical reference is Macdonald~\cite{Macdonald}; a textbook treatment
is \cite[Chapter~4]{BeckRobins}; and a primary-source proof of the more
general rational-cone statement, from which the polytope case follows, is in
Beck--Develin~\cite{BeckDevelin}. All three are retained because they are the
historical, textbook and cone-reciprocity citations respectively; the three
drafts merged here each used a different one. No assertion about a general
multivariate reciprocity law, about changing normal fans, or about a weighted
$q$-Ehrhart theorem is required anywhere below.

\section{Independent capacities and their reciprocity}
\label{sec:weighted}

This section is shared by both proofs, and it is the only place where they
share more than the two imports. Proof~A needs it only at $v=0$; Proof~B needs
the auxiliary capacity $v$ to stay alive until after reciprocity has been
applied, and then kills the non-maximal layers by setting it to zero. Setting
$v=0$ too early would hide exactly that operation.

\begin{proposition}[The independent-capacity polynomial]
\label{prop:count}
Let $\sigma$ be a preorder on $E$. For all nonnegative integers $u_i$ and $v$,
the number of lattice points of $K_\sigma(\mathbf u,v)$ equals the evaluation
at $(\mathbf u,v)$ of the polynomial
\begin{equation}
 \Ecal_\sigma(\mathbf u,v)=
 \sum_{a\in P_{\sigma^*}}
 \binom{v+n-|a|}{n-|a|}
 \prod_{i\in E}f_{a_i}(u_i),
 \label{eq:E}
\end{equation}
whose total degree is at most $n$.
\end{proposition}

\begin{proof}
Adjoin a new coordinate $0$ and put $I_0=E\cup\{0\}$ and
$I_i=U_\sigma(i)\cup\{0\}$ for $i\in E$. Lift the right-hand side of
\eqref{eq:minkowski} to
\[
 \widehat K=v\,\Delta_{I_0}+\sum_{i\in E}u_i\Delta_{I_i}
 \subseteq\R^{\{0\}\sqcup E}.
\]
It lies in the affine hyperplane $x_0+\sum_{i\in E}x_i=v+\sum_iu_i$, so
projection forgetting $x_0$ is an affine lattice isomorphism from that
hyperplane onto $\R^E$ and carries $\widehat K$ onto $K_\sigma(\mathbf u,v)$.

Apply Imported theorem~\ref{imp:lattice} with $F=\{0\}\sqcup E$, so $d=n$. Any
condition in \eqref{eq:draconian-conditions} involving the index $0$ is
automatic, because $I_0$ is the full ground set. For nonempty $A\subseteq E$
the remaining condition is
\begin{equation}
 a(A)\le|U_\sigma(A)|,\qquad U_\sigma(A)=\bigcup_{i\in A}U_\sigma(i).
 \label{eq:draconian-filter}
\end{equation}
If $A$ is a filter this says $a(A)\le|A|$; conversely those inequalities for
all filters imply \eqref{eq:draconian-filter}, since
$a(A)\le a(U_\sigma(A))\le|U_\sigma(A)|$. Filters of $\sigma$ are ideals of
$\sigma^*$, so the allowed $a$ are exactly the points of $P_{\sigma^*}$, and
the remaining coordinate $a_0=n-|a|$ is determined. Substituting into
\eqref{eq:postnikov} gives \eqref{eq:E}, the distinguished factor being
$f_{a_0}(v+1)=\binom{v+a_0}{a_0}$. The degrees of the factors in each summand
sum to $n$.

This filter identification of draconian vectors is the same one as
\cite[Lemma~4.1]{AC}; it is reproduced here to fix the dual orientation.
\end{proof}

\begin{proposition}[A weighted dual lattice-count polynomial]
\label{prop:weighted}
Put $L_{\tau^*}(\mathbf t)=\Ecal_{\tau^*}(\mathbf t,0)$. Then
\begin{equation}
 \boxed{\displaystyle
 L_{\tau^*}(\mathbf t)=\sum_{a\in P_\tau}\prod_{i\in E} f_{a_i}(t_i)}
 \qquad\text{and}\qquad
 L_{\tau^*}(\mathbf t)=\#\bigl(Q_{\tau^*}(\mathbf t)\cap\Z^E\bigr)
 \ \ (\mathbf t\in\N^E).
 \label{eq:L-formula}
\end{equation}
\end{proposition}

\begin{proof}
Take $\sigma=\tau^*$ in Proposition~\ref{prop:count}, so $\sigma^*=\tau$ and
$K_{\tau^*}(\mathbf t,0)=Q_{\tau^*}(\mathbf t)$. The distinguished factor is
$f_{a_0}(0+1)=1$ for every $a_0\ge0$, including zero weights; the remaining
factors are $f_{a_i}(t_i)$. The displayed formula also establishes
polynomiality directly.
\end{proof}

\begin{remark}[Overlap with the earlier package, and why the dual appears]
\label{rem:prior-enumerator}
Formula \eqref{eq:L-formula} is the $s=0$ specialization, with independent
capacities, of the rising-factorial enumerator of \cite{PriorReciprocity},
derived there in the same way from Postnikov's theorem and the same filter
identification of draconian vectors, both drafts crediting
\cite[Lemma~4.1]{AC}. It is re-derived here so that this document is
self-contained; it is not claimed as new. As for the orientation: in
$Q_{\tau^*}(\mathbf t)$ a source $i$ may send to precisely the $\tau$-lower
elements $D_\tau(i)$, and these are the neighbourhoods that turn the
draconian inequalities into the ideal inequalities defining $P_\tau$. Using
$U_\tau(i)$ at this point would interchange the preorder and its dual. This
orientation is checked once more in Section~\ref{sec:interior}.
\end{remark}

\begin{proposition}[Independent-capacity reciprocity]
\label{prop:reciprocity}
The polynomials of \eqref{eq:E} satisfy
\begin{equation}
 \Ecal_\sigma(-\mathbf u,-v)=(-1)^n\,\Ecal_\sigma(\mathbf u-\one,v-1)
 \label{eq:reciprocity}
\end{equation}
as an identity in all $n+1$ independent variables.
\end{proposition}

\begin{proof}
First take all $u_i$ and $v$ to be positive integers. Then
$K_\sigma(\mathbf u,v)$ is full-dimensional by Lemma~\ref{lem:transport}, and
$sK_\sigma(\mathbf u,v)=K_\sigma(s\mathbf u,sv)$ for $s\ge0$ by homogeneity of
the inequalities, so $L_K(s)=\Ecal_\sigma(s\mathbf u,sv)$ as polynomials in
$s$.

A lattice point $x$ is interior precisely when
\[
 x_i\ge1\ (i\in E),\qquad x(I)\le u(I)+v-1\quad(\emptyset\ne I\text{ an ideal}).
\]
Indeed, all the defining inequalities must be strict at an interior point,
since equality in a valid inequality with nonzero normal prevents a full ball
around it; and conversely simultaneous strictness of these finitely many
inequalities gives an open neighbourhood inside the polytope, redundant
inequalities causing no difficulty. The translation $x=z+\one$ is a bijection
onto the integer points with
\[
 z\ge0,\qquad z(I)\le(\mathbf u-\one)(I)+(v-1),
\]
whose capacities are nonnegative. Now \eqref{eq:ehrhart-reciprocity} and
Proposition~\ref{prop:count} give \eqref{eq:reciprocity} for every positive
integer capacity vector. Both sides are multivariate polynomials, and equality
on a Cartesian product of infinite sets forces polynomial equality: fix all
but one variable, use the univariate root bound, and repeat.
\end{proof}

\begin{remark}[Overlap with the earlier package]
\label{rem:prior-reciprocity}
Independently weighted capacity reciprocity is exactly what the manifest
entry for \cite{PriorReciprocity} already records, and its multivariate
reciprocity theorem is \eqref{eq:reciprocity}. The proof is supplied here so
that Proof~B is independent of that unrefereed draft; the mechanism is
credited to it as prior work within this collection and is not claimed as a
contribution. What goes beyond it is the conversion of \eqref{eq:reciprocity}
into arbitrary local coefficient weights in Section~\ref{sec:transfer}, and
the full $q$-evaluation.
\end{remark}

\begin{corollary}[Negative capacities as interior counts]
\label{cor:L-reciprocity}
For every positive integral vector $\mathbf c$,
\begin{equation}
 L_{\tau^*}(-\mathbf c)
 =(-1)^n\#\bigl(\Int Q_{\tau^*}(\mathbf c)\cap\Z^E\bigr).
 \label{eq:L-reciprocity}
\end{equation}
\end{corollary}

\begin{proof}
Fix $\mathbf c>0$ integral. By homogeneity $Q_{\tau^*}(s\mathbf c)
=sQ_{\tau^*}(\mathbf c)$ for $s\ge0$, and $Q_{\tau^*}(\mathbf c)$ is a
full-dimensional lattice polytope by Lemma~\ref{lem:transport}. Hence
$s\mapsto L_{\tau^*}(s\mathbf c)$ is its Ehrhart polynomial. Evaluate this
univariate polynomial at $s=-1$ and apply
\eqref{eq:ehrhart-reciprocity}.
\end{proof}

\section{Proof A, Engine 1: the support--interior identity}
\label{sec:interior}

This is the central combinatorial step of the first geometric engine. It works
with \emph{cumulative} supports. For $T\subseteq E$, put
\begin{align}
 B_\tau(T)&=\#\{a\in\Max(P_\tau):\supp(a)\subseteq T\},
 \label{eq:contained-B}\\
 \Lambda_\tau(T)&=\#\bigl(\Int Q_{\tau^*}(\one+\one_T)\cap\Z^E\bigr).
 \label{eq:interior-N}
\end{align}
The capacity vector $\one+\one_T$ is the distinctive idea of this route: a
capacity of $1$ everywhere, raised to $2$ exactly on $T$.

\begin{theorem}[Restricted supports as interior lattice points]
\label{thm:interior}
For every $T\subseteq E$,
\begin{equation}
 \boxed{B_\tau(T)=\Lambda_\tau(T).}
 \label{eq:interior-identity}
\end{equation}
In particular,
\begin{equation}
 B_\tau(T)=(-1)^nL_{\tau^*}(-\one-\one_T).
 \label{eq:restricted-negative}
\end{equation}
\end{theorem}

\begin{proof}
Assume first $n>0$. If $T=\emptyset$, then $B_\tau(T)=0$, since a maximal
point has total $n$; and an interior integral point of $Q_{\tau^*}(\one)$
would have all coordinates at least one and total strictly less than $n$,
which is impossible, so $\Lambda_\tau(T)=0$ as well.

Let $T\ne\emptyset$. Form a bipartite graph $G_T$ with left part $T$ and right
part $E\sqcup\{0\}$. Put an edge from $t$ to $0$ for every $t\in T$, and an
edge from $t$ to $i\in E$ exactly when
\begin{equation}
 i\le_\tau t.
 \label{eq:GT-edge}
\end{equation}
For nonempty $J\subseteq T$ we then have $N_{G_T}(J)=\{0\}\cup D_\tau(J)$.

\smallskip\noindent\emph{Left-hand draconian sequences.}
A vector $a\in\Dr_T(G_T)$ has total $n$ and satisfies
\begin{equation}
 a(J)\le |D_\tau(J)|\qquad(\emptyset\ne J\subseteq T).
 \label{eq:left-drac}
\end{equation}
Extend it by zero outside $T$. If $a\in P_\tau$, then \eqref{eq:left-drac}
follows from the ideal inequalities. Conversely, for any ideal $I$ let
$J=T\cap I$. If $J=\emptyset$ then $a(I)=0$; otherwise $D_\tau(J)\subseteq I$
and $a(I)=a(J)\le|D_\tau(J)|\le|I|$. Thus \eqref{eq:left-drac} is equivalent
to membership in $P_\tau$ for vectors of total $n$ supported in $T$, and
Lemma~\ref{lem:max} gives
\begin{equation}
 |\Dr_T(G_T)|=B_\tau(T).
 \label{eq:left-identification}
\end{equation}

\smallskip\noindent\emph{Right-hand draconian sequences.}
Write such a sequence as $(d_0,y)$ with $y\in\N^E$; its total is $|T|-1$. For
nonempty $S\subseteq E$ its condition is
\begin{equation}
 y(S)\le |T\cap U_\tau(S)|-1.
 \label{eq:right-subsets}
\end{equation}
Every subset containing the right vertex $0$ has all of $T$ as its neighbour
set, so those inequalities are automatic from the total and nonnegativity.
Conditions \eqref{eq:right-subsets} are equivalent to
\begin{equation}
 y(I)\le |T\cap I|-1
 \qquad\text{for every nonempty $\tau$-filter }I:
 \label{eq:right-filters}
\end{equation}
for one direction take $S=I$, noting $U_\tau(I)=I$; for the other apply
\eqref{eq:right-filters} to $U_\tau(S)$ and use $y(S)\le y(U_\tau(S))$. Taking
$I=E$ gives $|y|\le|T|-1$, so $d_0=|T|-1-|y|$ is determined and nonnegative.

Since the ideals of $\tau^*$ are the filters of $\tau$, the integral interior
points of $Q_{\tau^*}(\one+\one_T)$ are exactly the integer vectors $x$ with
\[
 x_i>0\ (i\in E),\qquad
 x(I)<|I|+|T\cap I|\quad(\emptyset\ne I\text{ a $\tau$-filter}).
\]
Full dimensionality was proved in Lemma~\ref{lem:transport}; consequently
strictness of these finitely many defining inequalities describes the ordinary
interior, even when some of them are redundant. Substituting $x=y+\one$ turns
these integral strict inequalities precisely into \eqref{eq:right-filters}. We
have therefore obtained a bijection $(d_0,y)\leftrightarrow y+\one$ between
right-hand draconian sequences and interior points, giving
\begin{equation}
 |\Dr_{E\sqcup\{0\}}(G_T)|=\Lambda_\tau(T).
 \label{eq:right-identification}
\end{equation}

\smallskip\noindent\emph{Isolated vertices and the duality step.}
If $D_\tau(T)\ne E$, there is an isolated right vertex. In that case the left
condition with $J=T$ would force $n\le|D_\tau(T)|<n$, and the right condition
for the singleton isolated vertex would force a nonnegative entry to be at
most $-1$; both draconian sets are empty. If $D_\tau(T)=E$, every right vertex
is adjacent to some left vertex, and the common neighbour $0$ connects all
left vertices, so $G_T$ is connected and Imported theorem~\ref{imp:duality}
applies. In both cases \eqref{eq:left-identification} and
\eqref{eq:right-identification} yield \eqref{eq:interior-identity}.

For $n=0$ use the usual zero-dimensional conventions: the only lattice point
is $\zero$, it is maximal and interior in $\R^0$, and both counts are one.
Finally, \eqref{eq:restricted-negative} follows from
Corollary~\ref{cor:L-reciprocity} with $\mathbf c=\one+\one_T$.
\end{proof}

The added vertex $0$ serves two distinct and indispensable purposes. On the
left it raises the total to $n$ and cancels the ``minus one'' in the
neighbourhood inequalities. On the right its entry stores the slack
$|T|-1-|y|$, and the ordinary interior then produces exactly the strictness
left over in the right neighbourhood inequalities.

\section{Proof A, completion: from negative capacities to exact supports}
\label{sec:completion}

For a finite coordinate downset $P\subseteq\N^E$, define the multiaffine
polynomial
\begin{equation}
 R_P(\mathbf z)=
 \sum_{\substack{B,S\subseteq E\\\one_B+\one_S\in P}}
 (-1)^{|B|+|S|}\zt B
 =\sum_{B\subseteq E}(-1)^{|B|}\eta_P(B)\,\zt B .
 \label{eq:R}
\end{equation}
The task is to prove $R_{P_\tau}=(-1)^nH_\tau$; by
Proposition~\ref{prop:local} this is Theorem~\ref{thm:main}.

\begin{lemma}[Two-value interpolation]
\label{lem:interpolation}
For the weighted polynomial of \eqref{eq:L-formula},
\begin{equation}
 R_{P_\tau}(\mathbf z)
 =\sum_{T\subseteq E}\zt T\omitprod T\,
       L_{\tau^*}(-\one-\one_T).
 \label{eq:two-value}
\end{equation}
\end{lemma}

\begin{proof}
Insert \eqref{eq:L-formula} into the right-hand side and interchange the
finite sums. The contribution of a fixed $a\in P_\tau$ factorizes as
\begin{equation}
 \prod_{i\in E}\bigl((1-z_i)f_{a_i}(-1)+z_i f_{a_i}(-2)\bigr).
 \label{eq:local-interpolation}
\end{equation}
Directly from the polynomial definition \eqref{eq:binomial-polynomial} of
$f_k$:
\begin{center}
\begin{tabular}{@{}c c c c@{}}
\toprule
$k$ & $f_k(-1)$ & $f_k(-2)$ & $(1-z)f_k(-1)+zf_k(-2)$\\
\midrule
$0$ & $1$ & $1$ & $1$\\
$1$ & $-1$ & $-2$ & $-(1+z)$\\
$2$ & $0$ & $1$ & $z$\\
$k\ge3$ & $0$ & $0$ & $0$\\
\bottomrule
\end{tabular}
\end{center}
Thus only $a\in\{0,1,2\}^E$ survive. The three nonzero local factors are
exactly the factors obtained by summing over all decompositions
$a=\one_B+\one_S$ with sign $(-1)^{|B|+|S|}$ and weight $\zt B$. This proves
\eqref{eq:two-value}.
\end{proof}

\begin{remark}[Seam: the same table in binomial language]
\label{rem:seam-table}
The four-row table above is the same content as the two-capacity collapse of
Section~\ref{sec:two-capacity}, which arises inside Proof~B by combining the
binomial identity \eqref{eq:binomial} multilinearly over $u_i\in\{1,2\}$: the
local maximal-layer factor there becomes $1$ at $b=0$ and $t$ at $b>0$,
exactly as the fourth column here becomes $1,-(1+z),z,0$ before the geometric
step. The two derivations were found independently and in different
languages, and the coincidence is one of the two seams where the two proofs
touch. Neither is used in the other's proof.
\end{remark}

\begin{proof}[Proof A of Theorem~\ref{thm:main}]
By Lemma~\ref{lem:interpolation} and Theorem~\ref{thm:interior},
\begin{align*}
 R_{P_\tau}(\mathbf z)
 &=(-1)^n\sum_{T\subseteq E}\zt T\omitprod T\,B_\tau(T)\\
 &=(-1)^n\sum_{a\in\Max(P_\tau)}
   \sum_{T\supseteq\supp(a)}\zt T\omitprod T.
\end{align*}
For any fixed $B\subseteq E$,
\begin{equation}
 \sum_{T\supseteq B}\zt T\omitprod T
 =\zt B\prod_{i\notin B}\bigl(z_i+(1-z_i)\bigr)=\zt B.
 \label{eq:bernstein-collapse}
\end{equation}
Consequently $R_{P_\tau}=(-1)^nH_\tau$, which is \eqref{eq:main}. Taking the
coefficient of $\zt B$ proves \eqref{eq:exact-support}; the exponents $n+|B|$
and $n-|B|$ have the same parity.

Finally, Proposition~\ref{prop:local} identifies $\Zeta_q(P_\tau,-q^{-1})$
with $R_{P_\tau}(q^{-1},\ldots,q^{-1})$, and Lemma~\ref{lem:max} identifies
the number of lower covers with support size. This proves the full formula
\eqref{eq:target} as an identity in $\Q(q)$. Substituting
\eqref{eq:exact-support} into \eqref{eq:reduction} for a general normalized
height proves Theorem~\ref{thm:height-intro}.
\end{proof}

\subsection{The proof compressed into a chain of equalities}

Writing $w_T(\mathbf z)=\zt T\prod_{i\notin T}(1-z_i)$, Proof~A is
\begin{align*}
 R_{P_\tau}(\mathbf z)
 &=\sum_T w_T(\mathbf z)L_{\tau^*}(-\one-\one_T)
 &&\text{by the local binomial table},\\
 &=(-1)^n\sum_Tw_T(\mathbf z)\Lambda_\tau(T)
 &&\text{by Ehrhart reciprocity},\\
 &=(-1)^n\sum_Tw_T(\mathbf z)B_\tau(T)
 &&\text{by bipartite duality},\\
 &=(-1)^nH_\tau(\mathbf z)
 &&\text{by \eqref{eq:bernstein-collapse}}.
\end{align*}
The $q$-zeta statement is the diagonal specialization, after the independent
incidence calculation of Section~\ref{sec:incidence}. The argument is finite
throughout: no convergence, no analytic continuation in a numerical parameter,
and no exchange of infinite sums is involved.

\subsection{Ordinary reciprocity is only the total-count shadow}
\label{sec:shadow}

At $q=1$, positive evaluations specialize to ordinary multichain counts.
Polynomial interpolation at the distinct nodes $m=2,\ldots,n+2$ shows that
this specialization of the $q$-zeta polynomial is the ordinary zeta
polynomial; this is the route the shipped verifier implements. A cleaner
justification avoids any specialization of a rational function in $q$
altogether: the ordinary zeta polynomial is represented by
$e_0^{\mathsf T}Z^{m-1}\mathbf j$, and since $Z=I+N$ with $N$ nilpotent,
$Z^{m-1}=\sum_j\binom{m-1}{j}N^j$ is polynomial in $m$ and evaluates to
$Z^{-2}$ at $m=-1$. Either way,
\[
 \Zeta(P_\tau,-1)=(-1)^n\#\Max(P_\tau).
\]
Both accounts are recorded because the merge must not appear to prove this
case again, but must explain why it falls out safely: this is the case already
claimed in the manifest and proved in \cite{PriorReciprocity}, and it is not
presented here as a new result.

Conversely, knowledge of this single total count cannot recover
\eqref{eq:exact-support}. The additional information is retained by the $2^n$
choices of the capacities $\one+\one_T$, and the two-value interpolation
converts it into support coefficients.

\section{Proof A, Engine 2: a private-element Euler identity}
\label{sec:private}

Engine~1 evaluated the cumulative counts $B_\tau(T)$ and recovered exact
supports afterwards, at the cost of importing Postnikov's bipartite duality.
Engine~2 attacks the exact support $B$ directly. It costs two Hall lemmas --
the point at which transitivity is spent -- and a longer geometric argument,
and it buys back the import: it uses Imported theorem~\ref{imp:lattice} alone,
applied twice, and thereby \emph{proves}, in the case it needs, the duality
that Engine~1 imports. Either engine completes Proof~A on its own. Both are
recorded, as announced in Remark~\ref{rem:import-trade}.

\subsection{An inequality description for sums of truncated simplices}

\begin{lemma}
\label{lem:inequalities}
Let $A_u\subseteq R$ for $u$ in a finite index set $U$, where $R$ is a finite
coordinate set. Then
\begin{equation}
 \sum_{u\in U}\Dz{A_u}
 =\left\{z\in\R_{\ge0}^{R}:
 z(J)\le\#\{u:A_u\cap J\ne\emptyset\}\ \text{ for all }J\subseteq R
 \right\}.
 \label{eq:rank-polytope}
\end{equation}
\end{lemma}

\begin{proof}
The inclusion from left to right holds separately for each summand. For the
reverse inclusion let $z$ satisfy the inequalities and let $w\in\R^R$.
Replacing the negative entries of $w$ by zero can only increase $w\cdot z$.
Let the distinct positive entries be $0<c_1<\cdots<c_k$, put $c_0=0$ and
$J_j=\{i:w_i\ge c_j\}$. Then
\[
 w\cdot z\le\sum_{j=1}^{k}(c_j-c_{j-1})z(J_j)
 \le\sum_{u\in U}\max\bigl(0,\max_{i\in A_u}w_i\bigr),
\]
and the last expression is the support function of the Minkowski sum. It
bounds $w\cdot z$ for every $w$, which places $z$ in that compact convex set
by the separating-hyperplane characterization. Empty $A_u$ contribute zero to
the support function. This is the general truncated-simplex-sum statement, for
an arbitrary indexed family; Lemma~\ref{lem:transport} is the instance in
which the family comes from a preorder.
\end{proof}

\subsection{Private elements and their alternating count}

Let $H$ be \emph{any} bipartite graph with disjoint vertex classes $R$ and
$C$. Write $r=|R|$, $c=|C|$, and $N(i)\subseteq C$ for the neighbours of
$i\in R$. Give each resource $i\in R$ an additional \emph{private} ground
element $p_i$, adjacent to resource $i$ alone; the augmented ground set is
\[
 U=\{p_i:i\in R\}\sqcup C .
\]
Let $\Ical_H$ consist of the subsets of $U$ that can be matched injectively
into $R$. Let $\Deg(H)$ be the set of distinct vectors
\begin{equation}
 d=(d_j)_{j\in C},\qquad
 d_j=\#\{i\in R:f(i)=j\},\qquad f(i)\in N(i),
 \label{eq:degree-vectors}
\end{equation}
that is, the set of degree vectors of full allocations -- not the set of
allocations themselves.

\begin{lemma}[Private-element Euler identity]
\label{lem:euler}
For every bipartite graph $H$ as above,
\begin{equation}
 (-1)^r\sum_{T\in\Ical_H}(-1)^{|T|}=|\Deg(H)|.
 \label{eq:euler}
\end{equation}
\end{lemma}

\begin{proof}
If $R=\emptyset$, the only independent subset is empty and the unique full
allocation has zero degree vector, so both sides are $1$. If some $i\in R$ has
$N(i)=\emptyset$, resource $i$ can serve only $p_i$; toggling $p_i$ preserves
membership in $\Ical_H$ and reverses the sign, so the left side is zero, while
the right side is zero because a full allocation into $C$ is impossible. We
may therefore assume $r\ge1$ and all $N(i)$ nonempty; in particular $c\ge1$.

\smallskip\noindent\emph{Step 1: turn the alternating sum into an Ehrhart
value.} For a ground element $u\in U$ let $A_u\subseteq R$ be its allowed
resources, so $A_{p_i}=\{i\}$ and $A_j=\{i:j\in N(i)\}$ for $j\in C$. Set
\begin{equation}
 K=\sum_{u\in U}\Dz{A_u}=[0,1]^R+\sum_{j\in C}\Dz{A_j}.
 \label{eq:K}
\end{equation}
Introduce a coordinate $0$ and homogenize each $\Dz{A_u}$ to
$\Delta_{\{0\}\cup A_u}$ on $F=\{0\}\sqcup R$. For a fixed nonnegative integer
$s$, dropping the zeroth coordinate is an affine lattice bijection from the
corresponding sum of simplices onto $sK$, since the sum of all coordinates is
fixed at $s|U|$. Introduce one further, full, simplex on $F$ with weight zero.

In \eqref{eq:draconian-conditions} a nonempty $V\subseteq U$ now imposes
$a(V)\le|\bigcup_{u\in V}A_u|$; subsets containing the full-simplex index give
no additional condition, and its exponent is $a_0=r-\sum_{u\in U}a_u$.
Consequently Imported theorem~\ref{imp:lattice} gives the polynomial identity
\begin{equation}
 \begin{gathered}
 L_K(s)=\sum_{a\in\mathcal A}\prod_{u\in U}f_{a_u}(s),\\
 \mathcal A=\Bigl\{a\in\N^U:a(V)\le\bigl|\textstyle\bigcup_{u\in V}A_u\bigr|
                           \text{ for all }V\subseteq U\Bigr\},
 \end{gathered}
 \label{eq:K-ehrhart}
\end{equation}
the sum being finite because $\sum_ua_u\le r$. The equality holds for all
nonnegative integers $s$ and is therefore an equality of polynomials. At
$s=-1$ a factor $f_{a_u}(-1)$ is $1$ for $a_u=0$, $-1$ for $a_u=1$, and $0$
for $a_u\ge2$; Hall's theorem identifies the surviving Boolean vectors with
$\Ical_H$. Thus
\begin{equation}
 L_K(-1)=\sum_{T\in\Ical_H}(-1)^{|T|}.
 \label{eq:alt-ehrhart}
\end{equation}

\smallskip\noindent\emph{Step 2: describe the interior lattice points.}
The cube in \eqref{eq:K} makes $K$ full-dimensional in $\R^R$. For
$J\subseteq R$ write $N(J)=\bigcup_{i\in J}N(i)$. Lemma~\ref{lem:inequalities}
gives
\[
 K=\{z\ge0:z(J)\le |J|+|N(J)|\ \text{ for all }J\subseteq R\}.
\]
Its interior is obtained by making the coordinate inequalities and all
nonempty-subset inequalities strict: an interior point cannot lie on a proper
supporting hyperplane, and conversely strict satisfaction of this finite
system gives a neighbourhood inside $K$. For an integer point $z$ put
$y=z-\one_R$. Then
\begin{equation}
 z\in\Int K\cap\Z^R
 \iff
 y\in\N^R\ \text{ and }\ y(J)\le|N(J)|-1
                   \quad(\emptyset\ne J\subseteq R).
 \label{eq:interior-y}
\end{equation}
The condition $J=R$ in particular bounds $\sum_iy_i$ by $c-1$. Note that the
shift removes the private cube and leaves \emph{strict} neighbourhood bounds,
that is $y(J)\le|N(J)|-1$ and not $y(J)\le|N(J)|$.

\smallskip\noindent\emph{Step 3: count the same vectors on the other side of
the graph.} Consider
\[
 P_H=\sum_{i\in R}\Delta_{N(i)}\subseteq\R^{C}.
\]
Apply Imported theorem~\ref{imp:lattice} once more, with a full simplex
$\Delta_C$ of weight zero and each displayed summand of weight one. Every
binomial factor then equals $1$. The resulting exponent vectors are precisely
the $y$ of \eqref{eq:interior-y}, completed by $a_0=c-1-\sum_iy_i$; subsets
containing the full-simplex index impose only the automatic bound
$a(J)\le c-1$. Hence
\begin{equation}
 \#(\Int K\cap\Z^R)=\#(P_H\cap\Z^{C}).
 \label{eq:interior-dual}
\end{equation}

Every degree vector in \eqref{eq:degree-vectors} is a lattice point of $P_H$.
Conversely, for $d\in P_H\cap\Z^{C}$ one has $\sum_jd_j=r$ and
$d(A)\le\#\{i:N(i)\cap A\ne\emptyset\}$ for all $A\subseteq C$; Hall's
theorem, applied to $d_j$ copies of each target $j$, provides an injective
assignment of these $r$ demands to the $r$ resources. All resources are used,
so this is a full allocation with degree vector $d$. Hence
$P_H\cap\Z^{C}=\Deg(H)$. Combining \eqref{eq:alt-ehrhart}, Ehrhart
reciprocity \eqref{eq:ehrhart-reciprocity} and \eqref{eq:interior-dual} proves
\eqref{eq:euler}.
\end{proof}

\begin{remark}[What the geometric argument does and does not assert]
The two sets in \eqref{eq:interior-dual} generally lie in different ambient
dimensions. The proof is an equality of their cardinalities obtained from
Imported theorem~\ref{imp:lattice}; it does not claim an affine bijection
between them, and it does not count matchings with multiplicity. Both
distinctions matter when the lemma is applied. Lemma~\ref{lem:euler} is
stated and proved for \emph{every} bipartite graph, not merely those arising
from preorders, so it is a reusable graph-theoretic statement in its own
right; the standalone test suite of Section~\ref{sec:computation} exercises it
in that generality.
\end{remark}

\subsection{Residual matchings and privatization}

Fix $B\subseteq E$ and put $R=E\setminus B$.

\begin{lemma}[Subtracting a Boolean vector]
\label{lem:residual}
For $y\in\N^E$, the vector $\one_B+y$ lies in $P_\tau$ exactly when the
demands specified by $y$ can be assigned injectively to resources in $R$,
using the relation $r\le_\tau i$.
\end{lemma}

\begin{proof}
The feasibility inequalities after subtraction are
\begin{equation}
 y(I)\le |I|-|I\cap B|=|I\cap R|\qquad\text{for every ideal }I.
 \label{eq:residual-ideal}
\end{equation}
For any set $A$ of demand types these imply
$y(A)\le y(D_\tau(A))\le|R\cap D_\tau(A)|$, and Hall's theorem gives the
matching. Conversely such a matching uses only resources in $R$, while the
base demands $\one_B$ can use their own resources in $B$; the combined
matching is feasible by Lemma~\ref{lem:allocation}.
\end{proof}

The residual graph still has targets in $R$ as well as targets in $B$: a
resource $r$ can serve a different target $u\in R$ when $r\le_\tau u$. The
next lemma removes those extra edges without changing any Boolean feasibility
test. It is the point at which transitivity does its work in this engine, and
it has no counterpart in Engine~1 or in Proof~B.

\begin{lemma}[Privatizing the residual targets]
\label{lem:private}
For $U\subseteq R$ and $V\subseteq B$, the set of demands $U\sqcup V$ can be
matched to $R$ under $\le_\tau$ if and only if $V$ can be matched to
$R\setminus U$ under $\le_\tau$.
\end{lemma}

\begin{proof}
The reverse implication follows by assigning each $u\in U$ to its own
resource. For the forward implication let $A\subseteq V$ and set
$I=D_\tau(A)$, an ideal by transitivity. Every matched demand in
$(U\cup V)\cap I$ must use a resource in $R\cap I$, so
\[
 |A|+|U\cap I|\le|(U\cup V)\cap I|\le|R\cap I|,
 \qquad\text{whence}\qquad
 |A|\le|(R\setminus U)\cap D_\tau(A)| .
\]
This is Hall's condition for matching $V$ to $R\setminus U$. No antisymmetry
was used.
\end{proof}

Now form the bipartite graph $H_B$ with resource set $R=E\setminus B$, target
set $B$, and edges
\begin{equation}
 r\mathrel{H_B}s\quad\Longleftrightarrow\quad r\le_\tau s,
 \label{eq:H-S}
\end{equation}
and adjoin to it a private ground element for each resource, identified with
the corresponding element of $R$. Lemmas~\ref{lem:residual}
and~\ref{lem:private} show that
\begin{equation}
 \eta_\tau(B)=\sum_{T\in\Ical_{H_B}}(-1)^{|T|}.
 \label{eq:eta-independent}
\end{equation}

\begin{lemma}[Maxima with prescribed support]
\label{lem:max-degrees}
The map $b\mapsto(b-\one_B)|_B$ is a bijection from the maximal elements of
$P_\tau$ with support $B$ onto $\Deg(H_B)$.
\end{lemma}

\begin{proof}
For such a $b$, Lemma~\ref{lem:max} gives $|b|=n$. The nonnegative vector
$y=b-\one_B$ is supported in $B$ and has sum $n-|B|=|R|$; by
Lemma~\ref{lem:residual} its demands can be assigned to $R$, and equality of
the totals forces every resource to be used. Thus $y|_B$ is a full-allocation
degree vector for $H_B$. Conversely a degree vector in $\Deg(H_B)$, extended
by zero outside $B$ and added to $\one_B$, is feasible, has sum $n$ and has
support exactly $B$, hence is maximal. The two constructions are inverse on
vectors; different allocations realizing the same degree vector produce the
same $b$, which is exactly what is required, since the final objects are
distinct degree vectors and not matchings counted with multiplicity.
\end{proof}

\begin{proof}[Proof A of Theorem~\ref{thm:main}, via Engine~2]
By \eqref{eq:eta-independent} and Lemma~\ref{lem:euler},
\[
 (-1)^{n-|B|}\eta_\tau(B)=|\Deg(H_B)|=N_\tau(B),
\]
the last equality being Lemma~\ref{lem:max-degrees}. This is
\eqref{eq:exact-support}, and it includes $B=\emptyset$, $B=E$ and
$E=\emptyset$ through the boundary cases of Lemma~\ref{lem:euler}. Summing
$(-1)^{|B|}\eta_\tau(B)\zt B$ over $B$ gives \eqref{eq:main}, and substitution
into \eqref{eq:reduction} proves \eqref{eq:target} and
Theorem~\ref{thm:height-intro} as before.
\end{proof}

\begin{remark}[Seam: one family, two descriptions]
\label{rem:seam-matroid}
The family $\Ical_{H_B}$ just used is precisely the family
$\{S\subseteq E:\one_B+\one_S\in P_\tau\}$ that Engine~1 meets in
\eqref{eq:exact-support}, and Proposition~\ref{prop:matroid} identifies it as
the independent-set family of a transversal matroid. In Engine~2 that
identification is load-bearing and is proved by Hall's theorem before
anything geometric happens; in Engine~1 it appears afterwards, as an
interpretation. This is the second of the two seams announced in
Section~\ref{sec:roadmap}.
\end{remark}

\section{Proof B: an involution transferring local weights}
\label{sec:transfer}

Proof~B is architecturally independent of Proof~A end to end. It uses no
incidence matrix, no M\"obius function, no Hall theorem, no bipartite graph,
no second polytope and no interior-point count. Its only inputs are
Section~\ref{sec:weighted} -- itself resting on the same two classical
theorems -- and the elementary structure of $P_\tau$. Keeping it is not
redundancy: it cross-validates an unrefereed result, and it is the only route
that reaches $\qint{-r}$ for every $r$.

\subsection{The \texorpdfstring{$v=0$}{v=0} pivot}

\begin{proposition}[Maximal-layer binomial identity]
\label{prop:binomial}
For independent indeterminates $u_i$,
\begin{equation}
 \sum_{a\in P_\tau}\prod_{i\in E}\binom{a_i-u_i-1}{a_i}
 =(-1)^n\sum_{b\in \Max(P_\tau)}\prod_{i\in E}\binom{b_i+u_i-2}{b_i}.
 \label{eq:binomial}
\end{equation}
\end{proposition}

\begin{proof}
Apply Proposition~\ref{prop:reciprocity} with $\sigma=\tau^*$, so that the
index set of \eqref{eq:E} is $P_{\sigma^*}=P_\tau$, and then set $v=0$. On the
left of \eqref{eq:reciprocity} the first factor of \eqref{eq:E} is
$\binom{n-|a|}{n-|a|}=1$. On the right it becomes
\[
 \binom{n-|a|-1}{n-|a|}=
 \begin{cases}1,&|a|=n,\\0,&|a|<n,\end{cases}
\]
which annihilates every layer below the top. This is a polynomial
specialization, legitimate because reciprocity was first established in all
$n+1$ variables. Lemma~\ref{lem:max} identifies the surviving layer with
$\Max(P_\tau)$. The remaining factors are exactly those of
\eqref{eq:binomial}, since $f_a(-u)=\binom{a-u-1}{a}$ and
$f_b(u-1)=\binom{b+u-2}{b}$ directly from \eqref{eq:binomial-polynomial}.
\end{proof}

The source of the top-layer restriction is now transparent, and it is worth
stating what it is not. The auxiliary capacity changes from $0$ to $-1$,
annihilating every positive slack $n-|a|$. There is \emph{no} geometric object
``of capacity $-1$'' anywhere in this argument, and no duality is invoked: the
selection of the maximal layer is a pure algebraic specialization. This is the
sharpest contrast with Proof~A, where the same selection is performed by a
bijection with interior lattice points (Engine~1) or with full allocations in
a residual graph (Engine~2).

\subsection{The involution and the transfer theorem}

For a formal power series over a commutative $\Q$-algebra, define
\begin{equation}
 (\Tr f)(z)=(1-z)f\left(\frac{-z}{1-z}\right).
 \label{eq:T}
\end{equation}
The substituted series has zero constant term, so the composition is
well-defined. The substitution $h(z)=-z/(1-z)$ satisfies $h(h(z))=z$ and
$1-h(z)=1/(1-z)$; hence $\Tr^2f=f$.

\begin{theorem}[Coefficient transfer]
\label{thm:transfer}
For any family $(f_i)_{i\in E}$ of formal power series over a commutative
$\Q$-algebra,
\begin{equation}
 \boxed{\displaystyle
 \sum_{a\in P_\tau}\prod_{i\in E}[z^{a_i}]f_i(z)
 =(-1)^n\sum_{b\in \Max(P_\tau)}\prod_{i\in E}[z^{b_i}](\Tr f_i)(z).}
 \label{eq:transfer}
\end{equation}
\end{theorem}

\begin{proof}
All indices on either side are at most $n$, and the first $n+1$ coefficients
of $\Tr f$ depend only on the first $n+1$ coefficients of $f$, so every $f_i$
may be truncated after degree $n$. The polynomials $(1-z)^m$, $0\le m\le n$,
form a basis of that polynomial space, and the identity is multilinear in the
$f_i$, so it suffices to take $f_i=(1-z)^{m_i}$ with $0\le m_i\le n$. For this
basis,
\[
 [z^a](1-z)^m=(-1)^a\binom ma=\binom{a-m-1}{a},
 \qquad \Tr\bigl((1-z)^m\bigr)=(1-z)^{1-m},
\]
\[
 [z^b](1-z)^{1-m}=\binom{b+m-2}{b},
\]
so that \eqref{eq:transfer} is \eqref{eq:binomial} with $u_i=m_i$.
Multilinearity and truncation finish the proof.
\end{proof}

Theorem~\ref{thm:transfer} is stated for arbitrary families of formal power
series over any commutative $\Q$-algebra, which is far more general than the
conjecture requires; that generality is what makes both
Theorem~\ref{thm:main} and Theorem~\ref{thm:allnegative} instances of one
statement.

\begin{remark}[A concrete coefficient rule]
\label{rem:Tcoeff}
If $f(z)=\sum_{k\ge0}f_kz^k$ and $\Tr f=\sum_{a\ge0}g_az^a$, then
\begin{equation}
 g_0=f_0,\qquad g_1=-f_0-f_1,\qquad
 g_a=\sum_{k=2}^a(-1)^k\binom{a-2}{k-2}f_k\quad(a\ge2),
 \label{eq:Tcoeff}
\end{equation}
by expanding $\Tr(z^k)=(-1)^kz^k(1-z)^{1-k}$. The transfer therefore uses a
triangular integral matrix, and is particularly simple to test without a
symbolic algebra system.
\end{remark}

\subsection{Proof B of the support identity}
\label{sec:proofB-support}

\begin{proof}[Proof B of Theorem~\ref{thm:main}]
Take $f_i(z)=(1-z)(1-z_iz)$, whose coefficient sequence is
$1,-(1+z_i),z_i,0,0,\ldots$, whereas
\[
 \Tr f_i(z)=1+\frac{z_iz}{1-z}
\]
has coefficient $1$ in degree zero and $z_i$ in every positive degree. The
left side of \eqref{eq:transfer} is then
\[
 \sum_{a\in P_\tau\cap\{0,1,2\}^E}
 \prod_{i:a_i=1}\bigl[-(1+z_i)\bigr]\prod_{i:a_i=2}z_i
 =R_{P_\tau}(\mathbf z),
\]
by grouping the pairs $(B,S)$ of \eqref{eq:R} according to
$a=\one_B+\one_S$, and the right side is
$(-1)^n\sum_{b}\prod_{i\in\supp(b)}z_i=(-1)^nH_\tau(\mathbf z)$. Hence
$R_{P_\tau}=(-1)^nH_\tau$, which is \eqref{eq:main}, and taking the
coefficient of $\zt B$ gives \eqref{eq:exact-support} in the form
\begin{equation}
 N_\tau(B)
 =(-1)^{n+|B|}\sum_{S\subseteq E}(-1)^{|S|}
       \mathbf 1_{\one_B+\one_S\in P_\tau}.
 \label{eq:fixedsupport}
\end{equation}
\end{proof}

This is a third, independent derivation of the exact-support identity: via
the transfer theorem, rather than via incidence inversion plus residual
matchings (Engine~2) or via cumulative counts plus two-value interpolation
(Engine~1). Since \eqref{eq:exact-support} is the strengthened statement on
which this whole document is built, all three routes to it are worth
recording.

Formula~\eqref{eq:fixedsupport} counts maximal elements with specified support
without enumerating their possibly large coordinate entries, and it displays a
nontrivial positivity statement for an alternating sum of feasibility
indicators. For $B=E$ the only feasible $S$ is the empty set, so the
coefficient of $\prod_iz_i$ is one, as expected from the unique all-positive
vector of sum $n$.

\subsection{The \texorpdfstring{$q$}{q}-zeta polynomial at every integer}
\label{sec:endpoint}

To complete Proof~B we must connect \eqref{eq:main} to $\Zeta_q$ without using
the incidence calculation of Section~\ref{sec:incidence}. This is done by
constructing the $q$-zeta polynomial outright. For $a\ge0$ and an integer $m$
use the rational-function convention
\begin{equation}
 \qbinom{m+a-2}{a}=\prod_{j=1}^a\frac{1-q^{m+j-2}}{1-q^j},
 \label{eq:gaussian}
\end{equation}
with empty product one; it extends the usual Gaussian binomial to the negative
upper arguments occurring below. The variable $q$ is indeterminate throughout.

\begin{proposition}[Gaussian endpoint formula]
\label{prop:endpoint}
For every integer $m$,
\begin{equation}
 \Zeta_q(P_\tau,\qint m)=\sum_{a\in P_\tau}q^{|a|}
                     \prod_{i\in E}\qbinom{m+a_i-2}{a_i}.
 \label{eq:endpoint}
\end{equation}
\end{proposition}

\begin{proof}
First let $m\ge2$ and condition a multichain on its final element $a$. Every
earlier point is coordinatewise at most $a$, hence lies in $P_\tau$ by the
downset property, and the choices in the separate coordinates are independent.
In a coordinate with endpoint $a_i$, the preceding $m-2$ entries form a weakly
increasing sequence in $\{0,\ldots,a_i\}$, whose weight enumerator is
$\qbinom{m+a_i-2}{a_i}$; the endpoint contributes $q^{a_i}$. For an elementary
verification of the local enumerator, let $R_{k,a}$ count weakly increasing
$k$-tuples with entries between zero and $a$ by the sum of their entries;
separating according to whether an entry $a$ occurs gives
$R_{k,a}=R_{k,a-1}+q^aR_{k-1,a}$ with $R_{0,a}=R_{k,0}=1$, and the Gaussian
polynomial $\qbinom{k+a}{a}$ satisfies the same recurrence and boundary
values.

To justify the remaining integers, regard $Q=q^m$ as a new variable. Each
local product in \eqref{eq:gaussian} is a polynomial in $Q$ of degree at most
$a_i$ over $\Q(q)$, so the right side of \eqref{eq:endpoint} is a polynomial
in $Q$ of degree at most $n$. Replacing $Q$ by $1+(q-1)X$ yields a polynomial
in $X$ whose values at $X=\qint m$, $m\ge2$, agree with \eqref{eq:qzeta-def}.
Such a polynomial is unique, because these arguments are infinitely many
distinct elements of $\Q(q)$. It is therefore the $q$-zeta polynomial, and the
same formula is valid at every integer $m$, including the negative ones.
\end{proof}

\begin{remark}[Two constructions at negative arguments]
\label{rem:two-constructions}
Proposition~\ref{prop:endpoint} and Lemma~\ref{lem:continuation} are two
different constructions of the $q$-zeta polynomial at negative arguments, and
neither implies the other. The endpoint formula is uniform in $m$ and is what
makes every $\qint{-r}$ reachable, through Lemma~\ref{lem:kernel} below; it is
special to coordinate downsets. The matrix continuation
$e_0^{\mathsf T}A^{m-1}\mathbf j$, with its rank-filtration diagonalization,
applies to an arbitrary finite poset with least element and an arbitrary
normalized height, and is what makes Theorem~\ref{thm:height-intro}
reachable. Both are kept.
\end{remark}

\begin{lemma}[The local negative kernel]
\label{lem:kernel}
For every $r\ge1$ and $a\ge0$,
\begin{equation}
 q^a\qbinom{a-r-2}{a}=[z^a]\prod_{j=0}^r(1-q^{-j}z).
 \label{eq:negativekernel}
\end{equation}
\end{lemma}

\begin{proof}
When $a\ge r+2$, the product \eqref{eq:gaussian} with $m=-r$ contains a factor
$1-q^0$, so both sides vanish. If $0\le a\le r+1$, rewrite the negative powers
by $1-q^{-k}=-q^{-k}(1-q^k)$ to get
\[
 q^a\qbinom{a-r-2}{a}=(-1)^aq^{-ra+\binom a2}\qbinom{r+1}{a},
\]
and expanding $\prod_{j=0}^r(1-q^{-j}z)$ by elementary symmetric functions
gives the same expression. This finite $q$-binomial identity also follows by
induction on $r$, multiplying by $1-q^{-r}z$.
\end{proof}

\begin{proof}[Proof B of \eqref{eq:target}]
By Proposition~\ref{prop:endpoint} and Lemma~\ref{lem:kernel} with $r=1$, the
left side of \eqref{eq:target} is the left side of \eqref{eq:transfer} with
every $f_i(z)=f(z)=(1-z)(1-q^{-1}z)$. Its transfer is
$\Tr f(z)=1+q^{-1}z/(1-z)$, whose coefficient is one in degree zero and
$q^{-1}$ in every positive degree. Hence Theorem~\ref{thm:transfer} gives
\[
 \Zeta_q(P_\tau,\qint{-1})
 =(-1)^n\sum_{b\in \Max(P_\tau)}\prod_{i:b_i>0}q^{-1}
 =(-1)^n\sum_{b\in\Max(P_\tau)}q^{-|\supp(b)|},
\]
and Lemma~\ref{lem:max} identifies that statistic with the $\cov$ of the
conjecture. No step of Section~\ref{sec:incidence} was used.
\end{proof}

\subsection{A shorter derivation once capacity reciprocity is known}
\label{sec:two-capacity}

The decisive interpolation can alternatively be compressed to the two capacity
values $u_i\in\{1,2\}$. Put $t=q^{-1}$. For every $a\ge0$,
\[
 [z^a](1-z)(1-tz)=(1-t)\binom{a-2}{a}+t\binom{a-3}{a},
\]
so multilinearity lets one combine \eqref{eq:binomial} over all choices
$u_i\in\{1,2\}$. On the maximal-layer side the corresponding local factor at
$b$ becomes
\[
 (1-t)\binom{b-1}{b}+t\binom bb
 =\begin{cases}1,&b=0,\\t,&b>0,\end{cases}
\]
which proves the conjecture directly. As Remark~\ref{rem:seam-table} records,
this two-capacity collapse is the same content as the two-value interpolation
table of Lemma~\ref{lem:interpolation}, in binomial rather than
$z$-polynomial language; the two were reached independently, from opposite
ends of the cluster of arguments merged here. The general transfer theorem is
kept nevertheless, because it is what exposes why the same method works at
every negative $q$-integer and not only at $\qint{-1}$.

\section{All negative evaluations: positivity and combinatorics}
\label{sec:negative}

Set $t=q^{-1}$ and $F_r(z)=\prod_{j=0}^r(1-t^jz)$. A direct substitution into
\eqref{eq:T} gives
\begin{equation}
 \Tr F_r(z)
 =\frac{\prod_{j=1}^r\bigl(1-(1-t^j)z\bigr)}{(1-z)^r}
 =\prod_{j=1}^r\left(1+\frac{t^jz}{1-z}\right).
 \label{eq:TF}
\end{equation}
To select a coefficient of $z^a$, choose the $s$ factors in which the
nonconstant summand is used; their contribution is
$e_s(t,t^2,\ldots,t^r)z^s(1-z)^{-s}$. For $a\ge1$ this gives
\[
 [z^a]\Tr F_r(z)=\sum_{s=1}^{\min(r,a)}\binom{a-1}{s-1}e_s(t,t^2,\ldots,t^r)
 =C_{r,a}(t),
\]
the degree-zero coefficient being one, in agreement with
\eqref{eq:Cdefinition}. One may also use the compact Gaussian form
\begin{equation}
 C_{r,a}(t)=\sum_{s=1}^{\min(r,a)}
 \binom{a-1}{s-1}t^{s(s+1)/2}\gausstbinom rs\qquad(a\ge1).
 \label{eq:Cgaussian}
\end{equation}

\begin{proof}[Proof of the identity and positivity in Theorem~\ref{thm:allnegative}]
Proposition~\ref{prop:endpoint} and Lemma~\ref{lem:kernel} express
$\Zeta_q(P_\tau,\qint{-r})$ as the left side of \eqref{eq:transfer} with every
$f_i=F_r$. Equation~\eqref{eq:TF} identifies the transferred coefficients as
$C_{r,a}(t)$, which proves \eqref{eq:allnegative}. Positivity and integrality
are immediate from \eqref{eq:Cdefinition}. The extreme-degree statements are
proved in Proposition~\ref{prop:degrees} below.
\end{proof}

Nothing in Proof~A goes past $\qint{-1}$; Theorem~\ref{thm:allnegative} is the
single largest piece of mathematics in this document that has only one proof.

\subsection{A colored-composition interpretation}

For $a>0$, an object counted by $C_{r,a}(t)$ consists of a nonempty set of
colors $1\le j_1<\cdots<j_s\le r$ together with a composition
$a=a_1+\cdots+a_s$ into $s$ positive parts; its weight is
$t^{j_1+\cdots+j_s}$. The colors are written in increasing order, so no factor
$s!$ occurs, and there are $\binom{a-1}{s-1}$ choices of composition. For
$a=0$ there is one empty object of weight one.

Consequently the signed negative evaluation counts pairs consisting of a
maximal lattice point $b$ and, independently at each coordinate, one such
colored composition of $b_i$. This explains the nonnegative coefficients
without any appeal to cancellation in the original expression for the $q$-zeta
polynomial. For instance,
\begin{align*}
 C_{1,a}(t)&=t &&(a\ge1),\\
 C_{2,a}(t)&=t+t^2+(a-1)t^3 &&(a\ge1),\\
 C_{3,a}(t)&=(t+t^2+t^3)+(a-1)(t^3+t^4+t^5)+\binom{a-1}{2}t^6 &&(a\ge1),
\end{align*}
the usual convention making the last binomial zero when $a=1,2$.

\subsection{The ordinary specialization at every \texorpdfstring{$r$}{r}}

Substituting $t=1$ in \eqref{eq:TF} gives $(1-z)^{-r}$, so
\begin{equation}
 C_{r,a}(1)=\binom{r+a-1}{a}.
 \label{eq:Cordinary}
\end{equation}
After the substitution $q^m=1+(q-1)X$, the $j$th local quotient in
\eqref{eq:gaussian} tends coefficientwise to $(X+j-2)/j$ as $q\to1$, so the
coefficientwise limit in the endpoint formula identifies the ordinary zeta
polynomial and yields
\begin{equation}
 (-1)^n \Zeta(P_\tau,-r)=
 \sum_{b\in \Max(P_\tau)}\prod_{i\in E}\binom{r+b_i-1}{b_i}.
 \label{eq:ordinary-r}
\end{equation}
At $r=1$ this is the unweighted identity already covered by
\cite{PriorReciprocity} and discussed in Section~\ref{sec:shadow}; it is not
claimed as new. At general $r$ it also follows directly by specializing all
$u_i=r+1$ in \eqref{eq:binomial}. The negative ordinary zeta values at $r>1$
appear nowhere else in this document.

\subsection{Sharp extreme degrees}
\label{sec:degrees}

\begin{lemma}[Minimum support size]
\label{lem:minsupp}
The smallest support size of an element of $\Max(P_\tau)$ is the number
$c_\tau$ of maximal equivalence classes of $\tau$.
\end{lemma}

\begin{proof}
A maximal equivalence class $F$ is a filter, so its complement is an ideal and
every $b\in\Max(P_\tau)$ satisfies $b(F)=n-b(E\setminus F)\ge|F|>0$. The
distinct maximal classes are disjoint, forcing at least one support coordinate
in each.

For the converse, choose a representative in each maximal class. For every
$i\in E$ choose a maximal class above the class of $i$ and send one unit from
$i$ to the chosen representative there; let $b_j$ count the units received by
$j$. Every unit assigned to a member of an ideal $I$ originates in $I$, since
its source lies below its destination, so $b(I)\le|I|$, and $|b|=n$; thus
$b\in\Max(P_\tau)$. Its support is exactly the set of chosen representatives,
because every maximal class receives at least the units from its own members.
\end{proof}

\begin{proposition}[Sharp extreme degrees]
\label{prop:degrees}
For $n,r\ge1$ let $H_{\tau,r}(t)=\sum_{b\in \Max(P_\tau)}\prod_iC_{r,b_i}(t)$.
Then $H_{\tau,r}$ has degree exactly $rn$ and leading coefficient one, its
smallest exponent is $c_\tau$, and the coefficient of $t^{c_\tau}$ is the
number of maximal points of support size $c_\tau$.
\end{proposition}

\begin{proof}
For every $a>0$ the smallest exponent in $C_{r,a}$ is one with coefficient
one, from the single color $1$, and positivity prevents cancellation; so the
smallest exponent in the summand indexed by $b$ is $|\supp(b)|$, and
Lemma~\ref{lem:minsupp} proves the assertion at the lower end.

A choice of $s\le a$ distinct colors has weight exponent at most $rs$, hence
at most $ra$. Equality with $ra$ is possible only when $a=1$: if $s<a$ the
first bound is strictly smaller than $ra$, and if $s=a\ge2$ then distinctness
of the colors makes their sum strictly smaller than $ra$. Thus
$\deg C_{r,a}\le ra$, with strict inequality for $a\ge2$, while
$C_{r,1}=t+\cdots+t^r$ is monic of degree $r$. Every maximal point has total
$n$, and the only one all of whose coordinates are at most one is $\one$; its
summand $(t+\cdots+t^r)^n$ is monic of degree $rn$, and every other summand
has smaller degree.
\end{proof}

This completes the proof of Theorem~\ref{thm:allnegative}.

\section{Consequences: matroids, feasible supports, and coefficients}
\label{sec:consequences}

\subsection{Two matroid invariants for the same coefficient}

Fix $B\subseteq E$. Make a bipartite graph whose sources are $E\setminus B$
and whose destinations are $E$, with an edge $i\to j$ whenever $i\le_\tau j$.
Let $M_B$ be the transversal matroid on the destination set $E$: a set $S$ is
independent when its elements can be matched to distinct sources. The matching
exchange property follows by taking alternating paths in the union of two
matchings, so these sets do form a matroid.

\begin{proposition}[Tutte evaluation]
\label{prop:matroid}
The independent sets of $M_B$ are exactly
\begin{equation}
 \Ical_B=\{S\subseteq E:\one_B+\one_S\in P_\tau\},
 \label{eq:matroid-independent}
\end{equation}
and $\rank M_B=n-|B|$. If $T_M(x,y)$ denotes the Tutte polynomial, then
\begin{equation}
 [\zt B]H_\tau(\mathbf z)=T_{M_B}(0,1).
 \label{eq:tutte}
\end{equation}
\end{proposition}

\begin{proof}
Membership in \eqref{eq:matroid-independent} is the family of inequalities
$|S\cap I|\le|I\setminus B|$ over ideals $I$. A matching implies them, because
every destination in $I$ must be matched from $I\setminus B$. Conversely, for
$A\subseteq S$ apply the inequality to $D_\tau(A)$:
\[
 |A|\le|S\cap D_\tau(A)|\le|D_\tau(A)\setminus B| ,
\]
which are Hall's conditions for matching $S$ into the available sources. There
are $n-|B|$ sources, and the destination set $E\setminus B$ can be matched
identically to them, which gives the rank.

By the subset expansion
$T_M(x,y)=\sum_{A\subseteq E}(x-1)^{\rank M-\rank(A)}(y-1)^{|A|-\rank(A)}$, at
$(x,y)=(0,1)$ only independent sets survive, giving
$T_M(0,1)=(-1)^{\rank M}\sum_{A\text{ independent}}(-1)^{|A|}$. Substitute
$\rank M_B=n-|B|$ and use \eqref{eq:exact-support}.
\end{proof}

The same coefficient carries a second, different standard invariant. By
Remark~\ref{rem:seam-matroid} the family $\Ical_B$ is the independent-set
family $\Ical_{H_B}$ of Engine~2, a transversal matroid of rank $n-|B|$ whose
private elements form a basis. Define the $h$-polynomial of its independence
complex directly by
\begin{equation}
 h_B(t)=\sum_{T\in\Ical_B}t^{|T|}(1-t)^{\,n-|B|-|T|} .
 \label{eq:h-polynomial}
\end{equation}
Its top coefficient is $(-1)^{n-|B|}\sum_T(-1)^{|T|}=N_\tau(B)$. So
\eqref{eq:tutte} and \eqref{eq:h-polynomial} are two standard matroid
invariants for one and the same number, reached from two different
literatures; \eqref{eq:h-polynomial} needs no matroid theorem at all, while
\eqref{eq:tutte} places the coefficient inside the Tutte formalism. Both are
recorded for that reason.

In either reading, the signs in the Boolean sum encode a familiar matroid
evaluation, but the equality with a count of maximal preorder points still
rests on Section~\ref{sec:interior} or Section~\ref{sec:private}.
Nonnegativity alone would not determine that count.

\subsection{Exactly which supports occur}
\label{sec:feasible}

Call $B\subseteq E$ \emph{cofinal} if every $r\in E$ satisfies $r\le_\tau s$
for some $s\in B$; equivalently, $B$ meets every maximal equivalence class of
$\tau$, that is, every maximal element of the quotient partial order. Two
statements are recorded, a qualitative one and a quantitative one; neither
follows from the other.

\begin{proposition}[Support feasibility, qualitative form]
\label{prop:support-feasible}
For every $B\subseteq E$, $N_\tau(B)>0$ if and only if $B$ is cofinal.
Consequently $\eta_\tau(B)=0$ on noncofinal supports, and on cofinal supports
its sign is exactly $(-1)^{n-|B|}$.
\end{proposition}

\begin{proof}
Let $C$ be a maximal class; its complement is an ideal, and a maximal point
has total $n$, so $a(E\setminus C)\le n-|C|$ forces $a(C)\ge|C|>0$, and the
support of $a$ meets $C$.

Conversely suppose $B$ meets every maximal class. Retain the unit of each
source $b\in B$ at its own destination; for each source $i\notin B$ choose a
maximal class above the class of $i$ and some $b\in B$ in that class, and send
the source's unit to $b$. This is allowed transport and uses every source, so
its destination vector has total $n$ and exact support $B$, hence is maximal.
The same conclusion follows from Lemma~\ref{lem:max-degrees}: positivity means
that every resource in $E\setminus B$ can choose some neighbour in $B$, the
choices not needing to be distinct because the degree vector records
multiplicities. The assertion about signs is Theorem~\ref{thm:main}.
\end{proof}

The quantitative counterpart is Lemma~\ref{lem:minsupp} together with
Proposition~\ref{prop:degrees}: the \emph{minimum} support size is the number
$c_\tau$ of maximal equivalence classes, and the coefficient at that exponent
counts the maximal points of that support size. Its constructive unit-routing
proof is kept because Proposition~\ref{prop:degrees}, and hence the degree
statement of Theorem~\ref{thm:allnegative}, depends on it.

It follows that the coefficient of $u^k$ in $H_\tau(u)$ is positive exactly
for $c_\tau\le k\le n$, with no internal zero coefficients. This does
\emph{not} assert log-concavity, real-rootedness, or any stronger coefficient
property.

\subsection{Coefficients of the support polynomial}

\begin{proposition}[Extreme and near-extreme coefficients]
\label{prop:top}
For $n\ge1$,
\begin{equation}
 [u^n]H_\tau(u)=1,\qquad
 [u^{n-1}]H_\tau(u)=\#\{(i,j)\in E^2:i\ne j,\ i\le_\tau j\}
 =\sum_{i\in E}\bigl(|U_\tau(i)|-1\bigr).
 \label{eq:top-coefficients}
\end{equation}
If the quotient preorder has exactly one maximal class $C$, then
$[u]H_\tau(u)=|C|$; if it has more than one maximal class, this coefficient is
zero. Equivalently, $[u]H_\tau(u)=\#\{s\in E:r\le_\tau s\text{ for all }r\in
E\}$. For $n\ge2$, writing $S_{ru}=E\setminus\{r,u\}$,
\begin{equation}
 [u^{n-2}]H_\tau(u)=
 \sum_{\{r,u\}\subseteq E}
 \left(
 |U_\tau(r)\cap S_{ru}|\,|U_\tau(u)\cap S_{ru}|
 -\binom{|U_\tau(r)\cap U_\tau(u)\cap S_{ru}|}{2}
 \right).
 \label{eq:next-coefficient}
\end{equation}
\end{proposition}

\begin{proof}
The only positive-coordinate integer vector of total $n$ is $\one$. A vector
of total $n$ with support of size $n-1$ is uniquely $\one-e_i+e_j$ with $i\ne
j$, and it satisfies the ideal inequalities exactly when every ideal
containing $j$ also contains $i$, which is equivalent to $i\le_\tau j$; for
necessity take the ideal $D_\tau(j)$. For singleton support the vector must be
$ne_j$, and the transport criterion says that every source lies below $j$,
which in a finite preorder happens exactly when there is a unique maximal
class and $j$ lies in it.

For \eqref{eq:next-coefficient}, use Lemma~\ref{lem:max-degrees}: fix the two
residual resources $r,u$, and let their target sets inside $S_{ru}$ have sizes
$a,b$ with intersection of size $d$. There are $ab$ ordered choices. Two
distinct choices produce the same degree vector exactly when two distinct
common targets are exchanged between $r$ and $u$, so each of the $\binom d2$
unordered pairs of common targets creates exactly one excess count, and no
degree vector has any other multiplicity. The number of distinct vectors is
therefore $ab-\binom d2$. Summing over residual pairs proves the formula.
\end{proof}

Formula \eqref{eq:next-coefficient} is the only third-from-the-top
coefficient formula available here, and the product rule below is the only
statement about disjoint unions; they came from different drafts and neither
subsumes the other. In particular, the negative $q$-evaluation has pole order
exactly $n$ at $q=0$:
\[
 (-1)^n\Zeta_q(P_\tau,-q^{-1})
 =q^{-n}+c'_\tau q^{-(n-1)}+\text{higher powers of }q,
\]
where $c'_\tau$ counts the ordered nonreflexive comparisons in
\eqref{eq:top-coefficients}. Equivalent but distinct elements contribute in
both directions.

\subsection{Products and an inclusion--exclusion algorithm}

If $\tau=\tau_1\sqcup\tau_2$ is a disjoint union on disjoint ground sets, then
$P_\tau=P_{\tau_1}\times P_{\tau_2}$, its maximal points are products of
maximal points, and therefore
\[
 H_\tau(\mathbf z)=H_{\tau_1}(\mathbf z_{E_1})H_{\tau_2}(\mathbf z_{E_2}),
\]
with no additional interaction between components.

The cumulative counts can be inverted on the Boolean lattice:
\begin{equation}
 [\zt B]H_\tau(\mathbf z)
 =\sum_{T\subseteq B}(-1)^{|B|-|T|}B_\tau(T)
 =(-1)^n\sum_{T\subseteq B}(-1)^{|B|-|T|}L_{\tau^*}(-\one-\one_T).
 \label{eq:support-inversion}
\end{equation}
Once the $2^n$ cumulative values are known, all exact-support counts follow in
$O(n2^n)$ additions by the Boolean M\"obius transform. This is an exact finite
procedure, not a claim that obtaining the cumulative counts is
polynomial-time.

\subsection{Arbitrary heights and a multivariate matrix identity}
\label{sec:heights}

Theorem~\ref{thm:height-intro} was proved in Section~\ref{sec:completion}. Two
consequences deserve separate statements.

\begin{corollary}[Linear weights]
\label{cor:height}
For positive integers $w_i$ and $h(a)=\sum_iw_ia_i$,
\begin{equation}
 \Zeta_{q,h}(P_\tau,\qint{-1})
   =(-1)^n\sum_{b\in\Max(P_\tau)}q^{-\sum_{i\in\supp(b)}w_i}.
 \label{eq:linear-height}
\end{equation}
\end{corollary}

\begin{proof}
Substitute $h(\one_{\supp(b)})=\sum_{i\in\supp(b)}w_i$ in
\eqref{eq:height-intro}. The exponent involves the Boolean support vector and
not $h(b)$; the M\"obius function in \eqref{eq:reduction} is exactly what
enforces that distinction.
\end{proof}

\begin{corollary}[Multivariate matrix identity]
\label{cor:multi}
Let $z_i$ be commuting indeterminates, put $\mathbf z^a=\prod_iz_i^{a_i}$ and
$D(\mathbf z)=\operatorname{diag}(\mathbf z^a)_{a\in P_\tau}$. Then, over
$\Q(z_i:i\in E)$,
\begin{equation}
 e_0^{\mathsf T}\bigl(ZD(\mathbf z)\bigr)^{-2}\mathbf j
    =(-1)^n\sum_{b\in\Max(P_\tau)}\prod_{i\in\supp(b)}z_i^{-1}.
 \label{eq:multi}
\end{equation}
\end{corollary}

\begin{proof}
The purely algebraic equality
$e_0^{\mathsf T}(ZD)^{-2}\mathbf j=e_0^{\mathsf T}\mu D^{-1}\mu\mathbf j$
still holds, and \eqref{eq:mobius} reduces it to
$\sum_B(-1)^{|B|}\eta_\tau(B)\prod_{i\in B}z_i^{-1}$. Apply
Theorem~\ref{thm:main}.
\end{proof}

Equation \eqref{eq:multi} is a multivariate \emph{matrix} identity; it is not
a claim that some multivariate interpolation polynomial has been constructed,
and no such object is needed. Setting all $z_i=1$ is legitimate because $Z$ is
invertible, and returns $\Zeta(P_\tau,-1)=(-1)^n\#\Max(P_\tau)$ as in
Section~\ref{sec:shadow}.

\section{Explicit examples}
\label{sec:examples}

\subsection{Antichains, discrete preorders, and a single equivalence class}

For an $n$-element antichain the defining inequalities reduce to $0\le x_i\le
1$, so $P_\tau=\{0,1\}^E$ and the only maximal point is $\one$:
\[
 H_\tau(\mathbf z)=\prod_iz_i,\qquad
 \Zeta_q(P_\tau,-q^{-1})=(-1)^nq^{-n},\qquad
 \Zeta_q(P_\tau,X)=X^n,
\]
the last because a multichain is determined coordinatewise by its number of
ones, which may be $0,\ldots,m-1$. At the other negative arguments,
$(-1)^n\Zeta_q(P_\tau,\qint{-r})=(t+t^2+\cdots+t^r)^n$. This example checks
both the sign and the shift in the length of the multichains.

At the opposite extreme let all $n$ elements be equivalent. The only nonempty
ideal is $E$, so $P_\tau=\{a\ge0:|a|\le n\}$ and the maximal points are the
weak compositions of $n$ into $n$ parts. For a prescribed support $B$ of size
$b\ge1$ the positive coordinates form a composition of $n$ into $b$ positive
parts, so, by placing $b-1$ separators in the $n-1$ internal gaps,
\begin{equation}
 [\zt B]H_\tau(\mathbf z)=\binom{n-1}{b-1},\qquad
 H_\tau(u)=\sum_{b=1}^n\binom nb\binom{n-1}{b-1}u^b.
 \label{eq:one-class}
\end{equation}
For $n=3$ this is $3u+6u^2+u^3$ with $10$ maximal points, and for $n=5$ it is
$5u+40u^2+60u^3+20u^4+u^5$ with $126$. The example also shows that an
equivalence class must not be collapsed to one coordinate: its cardinality
affects both the polytope and the support statistic.

\subsection{Chains}

For the partial order $1<2<\cdots<n$ the maximal vectors are described by
$a_1+\cdots+a_i\le i$ for $1\le i<n$ together with $|a|=n$. For $n=3$ they are
\[
 (0,0,3),\quad(0,1,2),\quad(0,2,1),\quad(1,0,2),\quad(1,1,1),
\]
so $H_\tau(u)=u+3u^2+u^3$, with full multivariate form
$z_3+2z_2z_3+z_1z_3+z_1z_2z_3$. The nonzero supportwise counts and their
signed Boolean sums are
\begin{center}
\begin{tabular}{@{}lrr@{}}
\toprule
$B$ & $N_\tau(B)$ & $\eta_\tau(B)$\\
\midrule
$\{3\}$ & 1 & 1\\
$\{1,3\}$ & 1 & $-1$\\
$\{2,3\}$ & 2 & $-2$\\
$\{1,2,3\}$ & 1 & 1\\
\bottomrule
\end{tabular}
\end{center}
All other supports have both counts zero, because they are not cofinal. For
$B=\{2,3\}$ the feasible $S$ are $\emptyset,\{1\},\{2\},\{3\}$, with
alternating sum $-2$; for $B=\{3\}$ the independent-subset counts by size in
the private residual graph are $1,3,3$, with alternating sum $1$. Hence
$\Zeta_q(P_\tau,-q^{-1})=-q^{-1}-3q^{-2}-q^{-3}$, and at the next negative
argument
\[
 (-1)^3\Zeta_q(P_\tau,\qint{-2})=t+4t^2+9t^3+9t^4+6t^5+t^6 .
\]
For the nonlinear height $h(a)=|a|^2$, Theorem~\ref{thm:height-intro} instead
gives
\[
 \Zeta_{q,h}(P_\tau,\qint{-1})=-q^{-1}-3q^{-4}-q^{-9},
\]
which is $-353/512$ at $q=2$, against $-11/8$ in the linear case. The verifier
obtains both by interpolation from positive multichain values, with a degree
bound of $9$ in the nonlinear case, without ever using a negative-evaluation
formula.

For $n=5$ direct enumeration gives $132$ lattice points, $42$ maximal vectors,
\[
 H_\tau(u)=u+10u^2+20u^3+10u^4+u^5,
 \qquad \Zeta_2(P_\tau,-1/2)=-\tfrac{197}{32}.
\]
This example has a unique maximum, so its least support size is one, and its
coefficient at $u^4$ is the ten ordered comparisons $i<j$.

\subsection{A fork, and why the direction of the preorder matters}

For $V$: $1<2$ and $1<3$, the three maximal points are $(0,1,2)$, $(0,2,1)$
and $(1,1,1)$, so
\[
 H_V(\mathbf z)=2z_2z_3+z_1z_2z_3,\qquad
 \Zeta_q(P_V,-q^{-1})=-2q^{-2}-q^{-3},\qquad H_V(t)=2t^2+t^3 .
\]
There are two maximal classes, which explains the absence of a $t$ term before
any computation. This small example already distinguishes maximal-point
support from rank: all three points have rank three, but their support sizes
are two, two and three, so replacing the cover statistic by the rank would
give an incorrect formula.

For the dual preorder $V^*$, relabelled so that $3$ lies above the
incomparable $1,2$, one gets instead $H_{V^*}(t)=t+2t^2+t^3$. Both
lattice-point posets have $12$ elements, but their numbers of maxima are $3$
and $4$. Neither the maximal-support polynomial nor its $q$-zeta
specialization should therefore be confused with the different
support-enumerator duality studied in~\cite{DaiEtAl}.

\subsection{A diamond}

For the partial order $1<2$, $1<3$, $2<4$, $3<4$ one finds $|P_\tau|=38$,
$|\Max(P_\tau)|=12$ and
\[
 H_\tau(t)=t+5t^2+5t^3+t^4 .
\]
This is the only four-element example worked out here.

\subsection{The five-element example in the source}

Take $E=\{a,b,c,d,e\}$ with $a\sim c$, with $a,c<e$, with $b<d$ and $b<e$, and
with no other comparisons except those forced by reflexivity and transitivity.
This is the running example of~\cite[Sections~3.1 and~4.4]{AC}. Besides
nonnegativity, its polytope is described by
\begin{align*}
 x_a+x_c&\le2,& x_b&\le1,\\
 x_a+x_b+x_c+x_e&\le4,& x_b+x_d&\le2,\\
 x_a+x_b+x_c+x_d+x_e&\le5.
\end{align*}
There are $92$ lattice points and $18$ maximal ones. The maximal classes are
$\{d\}$ and $\{e\}$, so every support contains both $d$ and $e$. The complete
exact-support distribution is
\begin{center}
\begin{tabular}{@{}l r@{\qquad}l r@{}}
\toprule
Support & Count & Support & Count\\
\midrule
$\{d,e\}$ & 2 & $\{a,b,d,e\}$ & 2\\
$\{a,d,e\}$ & 4 & $\{a,c,d,e\}$ & 2\\
$\{b,d,e\}$ & 1 & $\{b,c,d,e\}$ & 2\\
$\{c,d,e\}$ & 4 & $\{a,b,c,d,e\}$ & 1\\
\bottomrule
\end{tabular}
\end{center}
Thus
\begin{align}
 H_\tau(\mathbf z)
 &=z_dz_e\bigl(2+4z_a+z_b+4z_c
       +2z_az_b+2z_az_c+2z_bz_c+z_az_bz_c\bigr),
 \label{eq:source-multivariate}\\
 H_\tau(u)&=2u^2+9u^3+6u^4+u^5,
 \label{eq:source-diagonal}
\end{align}
and Theorem~\ref{thm:main} gives
\[
 \Zeta_q(P_\tau,-q^{-1})
 =-\frac{2q^3+9q^2+6q+1}{q^5}
 =-\frac{(2q+1)(q^2+4q+1)}{q^5},
\]
which agrees with the expression printed on page~15 of~\cite{AC}. The support
table is a refinement of that expression, not merely its total value $-18$ at
$q=1$; \eqref{eq:source-multivariate} is finer still, and is the multivariate
refinement of the source's own example. At the next negative argument,
Theorem~\ref{thm:allnegative} gives
\begin{align*}
 (-1)^5\Zeta_q(P_\tau,\qint{-2})
 ={}&2t^2+13t^3+41t^4+76t^5+94t^6\\
    &+76t^7+40t^8+11t^9+t^{10},
\end{align*}
whose lowest exponent $2$ agrees with the two maximal equivalence classes, as
Proposition~\ref{prop:degrees} requires.

\begin{table}[htbp]
\centering
\small
\begin{tabular}{@{}lrrlr@{}}
\toprule
Preorder & $|P_\tau|$ & $|\Max(P_\tau)|$ & $H_\tau(t)$
 & $\Zeta_q(P_\tau,-q^{-1})|_{q=2}$\\
\midrule
Antichain, $n=3$ & 8 & 1 & $t^3$ & $-1/8$\\
Chain, $n=3$ & 14 & 5 & $t+3t^2+t^3$ & $-11/8$\\
$V$ (fork), $n=3$ & 12 & 3 & $2t^2+t^3$ & $-5/8$\\
$V^*$, $n=3$ & 12 & 4 & $t+2t^2+t^3$ & $-9/8$\\
One equivalence class, $n=3$ & 20 & 10 & $3t+6t^2+t^3$ & $-25/8$\\
Diamond, $n=4$ & 38 & 12 & $t+5t^2+5t^3+t^4$ & ---\\
Chain, $n=5$ & 132 & 42 & $t+10t^2+20t^3+10t^4+t^5$ & $-197/32$\\
One equivalence class, $n=5$ & 252 & 126 & $5t+40t^2+60t^3+20t^4+t^5$
 & $-681/32$\\
Source example, $n=5$ & 92 & 18 & $2t^2+9t^3+6t^4+t^5$ & $-65/32$\\
\bottomrule
\end{tabular}
\caption{Exact examples, all computed independently by the included
verifiers. The diamond has relations $1<2$, $1<3$, $2<4$, $3<4$. The last
column is the negative $q$-evaluation at $q=2$; the seven rows for which it is
given are also recorded in \texttt{data/examples.csv}.}
\label{tab:examples}
\end{table}

\section{Why the hypotheses cannot simply be discarded}
\label{sec:limits}

Three examples are recorded, carrying four distinct failure statements
between them. Each bounds a different hypothesis, and none of them can be
deduced from the others: the first shows that purity of the downset is not
enough, the second that transitivity is not implied by the rest of the
transport picture, and the third that unit capacities are essential.

\subsection{Purity alone is insufficient: a pure downset, two failures}
\label{sec:pure-downset}

It is natural to ask whether \eqref{eq:target}, or the transfer identity
\eqref{eq:transfer} behind Proof~B, holds for every finite pure coordinate
downset in $\N^E$ containing the Boolean cube. Neither does. Let
\begin{equation}
 P=[\zero,(1,1,1)]\cup[\zero,(3,0,0)]\subseteq\N^3,
 \label{eq:counterexample}
\end{equation}
a finite graded coordinate downset containing every Boolean vector, whose two
maximal elements $(1,1,1)$ and $(3,0,0)$ both have rank three.

\smallskip\noindent\emph{Failure of the supportwise identity and of the
$q$-zeta value.} For $B=\{1\}$ the terms with $S\subseteq\{2,3\}$ cancel,
while $S=\{1\}$ contributes $-1$; so $\eta_P(\{1\})=-1$, whereas
$(-1)^{3-|B|}N_P(\{1\})=+1$. The only other nonvanishing $\eta_P$ is at
$B=\{1,2,3\}$, where it is $1$, so \eqref{eq:reduction} gives
\[
 \Zeta_q(P,-q^{-1})=q^{-1}-q^{-3},
\]
whereas the proposed extension would give $-q^{-1}-q^{-3}$. At $q=2$ the
discrepancy is $3/8$ against $-5/8$.

\smallskip\noindent\emph{Failure of the transfer identity.} With
$f_1(z)=z^2$ and $f_2=f_3=1$, the two sides of \eqref{eq:transfer} are $1$ and
$-1$. This is immediate from $\Tr(z^2)=z^2/(1-z)$ and $\Tr(1)=1-z$: the left
side counts the single point $(2,0,0)\in P$, while on the right only
$b=(3,0,0)$ contributes, with weight $1$ and sign $(-1)^3$.

Both statements are kept. They are genuinely different: the first says the
conclusion of Theorem~\ref{thm:main} fails for $P$, the second says the
much more general machine of Theorem~\ref{thm:transfer} fails for $P$, so
capacity reciprocity is a real restriction inside Proof~B and not an artifact
of its presentation. This counterexample locates a genuine restriction of the
theorem: its proofs use the preorder allocation structure, not merely purity
and product-of-chains intervals.

\subsection{A reflexive, nontransitive incidence system}
\label{sec:nontransitive}

One could instead try to prove the theorem for any downset generated by
assigning unit sources to allowed destinations, with each source allowed to
stay in place. The following example rules that out, and it is not covered by
Section~\ref{sec:pure-downset}, because here the transport representation is
present and only transitivity is missing.

On three coordinates give the sources the neighbourhoods
\[
 A_1=\{1,2\},\qquad A_2=\{2,3\},\qquad A_3=\{1,3\}.
\]
The relation is reflexive but not transitive. Its distinct full-assignment
vectors are the six permutations of $(2,1,0)$ together with $(1,1,1)$. Let $P$
be their coordinate downset. All seven generators are maximal, and
\begin{equation}
 H_P(\mathbf z)=2(z_1z_2+z_1z_3+z_2z_3)+z_1z_2z_3.
 \label{eq:counter-support}
\end{equation}
The downset contains every vector of total at most two, and among total-three
vectors exactly the seven just listed. The general Boolean localization
formula \eqref{eq:local} gives
\begin{equation}
 -R_P(\mathbf z)
 =z_1+z_2+z_3+2(z_1z_2+z_1z_3+z_2z_3)+z_1z_2z_3 .
 \label{eq:counter-boolean}
\end{equation}
For instance, at the coefficient of $z_1$ the feasible $S$ are all the proper
subsets of $\{1,2,3\}$, whose alternating sum is $1-3+3=1$; yet there is no
maximal point supported only at coordinate one. In diagonal form the true
maximal-support polynomial is $6u^2+u^3$, whereas the signed negative
$q$-zeta value is $3u+6u^2+u^3$.

Reflexivity, purity, unit supplies, coordinate-box intervals and a transport
representation are therefore jointly insufficient. Transitivity enters
decisively when arbitrary neighbourhood constraints are replaced by ideal or
filter inequalities, which is exactly what Lemma~\ref{lem:private} and the
filter identification inside Proposition~\ref{prop:count} exploit.

\subsection{Arbitrary positive capacities do not satisfy the same formula}

The auxiliary polynomial $L_{\tau^*}(\mathbf t)$ has independent capacity
variables, but this does not extend the main $q$-formula unchanged to
$Q_\tau(\mathbf c)$ with arbitrary capacities. Take one coordinate of capacity
two: the lattice-point poset is the chain $0<1<2$, and \eqref{eq:ternary}
gives
\[
 \Zeta_q(P,-q^{-1})=1-(1+q^{-1})+q^{-1}=0 .
\]
Its sole maximal point has support size one, so neither a sign based on
ambient dimension nor one based on maximal rank can produce the asserted
support formula. The unit-capacity model in the main theorem is essential.

\subsection{What the result does not settle}
\label{sec:notsettled}

This document does not establish the other real-rootedness,
magic-positivity or $\gamma$-positivity conjectures for preorder polytopes,
and nonnegativity of the coefficients in \eqref{eq:allnegative} does not imply
real-rootedness or log-concavity; no such inference is made. As recorded in
Section~\ref{sec:nonduplication}, $H_\tau$ counts only maximal points, whereas
the preorder $h$-polynomial of \cite[Section~5]{AC} counts support sizes over
all lattice points; these are different polynomials.

We also do not provide a canonical pointwise bijection between $B_\tau(T)$ and
$\Lambda_\tau(T)$: in Engine~1 that equality uses the imported bipartite
duality, and a triangulation-dependent bijection could be sought through that
theorem but is not needed. Similarly, no direct sign-reversing involution
realizing \eqref{eq:exact-support} is constructed. In Proof~B the operator
$\Tr$ is an algebraic involution on local series, not a combinatorial matching
of individual global summands; the colored-composition model of
Section~\ref{sec:negative} gives a positive interpretation of the final answer
without pretending to supply such a matching.

Two open problems are recorded rather than solved. First, characterize the
finite coordinate downsets for which the transfer identity
\eqref{eq:transfer} holds for every family of local series; the example of
Section~\ref{sec:pure-downset} shows the class is strictly smaller than the
pure downsets containing the Boolean cube. Second, characterize the finite
downsets satisfying the supportwise identity \eqref{eq:exact-support}. Neither
is claimed as an additional solved problem here.

\section{Exact computation and reproducibility}
\label{sec:computation}

Three independent verifiers are shipped, one from each of the merged
packages. They test three different theorems and none subsumes the others, so
all three are retained rather than unified:
\begin{itemize}
\item \texttt{code/verify.py} --- the support/interior/transport suite, with
 four random-construction modes and the weighted lattice polynomial at both
 nonnegative and negative capacities;
\item \texttt{code/verify\_transfer.py} --- the $\qint{-r}$ polynomial
 identities at $r=1,\ldots,4$, arbitrary signed local-kernel transfer tests,
 direct capacity and reciprocity counts, and multichain dynamic-programming
 interpolation;
\item \texttt{code/verify\_supportwise.py} --- the inverse-matrix evaluations,
 the nonlinear-height interpolation, and the standalone bipartite-lemma suite;
\item \texttt{code/test\_unit.py} --- ten fast normalization and boundary
 tests.
\end{itemize}
All of them use only the Python standard library, all counts use integers, and
all interpolations and polynomial evaluations use
\texttt{fractions.Fraction}. There is no floating-point tolerance, computer
algebra service or hidden dataset; the only floating-point quantity recorded
anywhere is elapsed wall-clock time. A failed equality raises an exception
instead of silently recording approximate agreement.

Note that the runs used \emph{different seeds} and \emph{different sampling
depths}, and this has been preserved deliberately: \texttt{verify.py} and
\texttt{verify\_transfer.py} use seed \texttt{20260920}, while
\texttt{verify\_supportwise.py} uses seed \texttt{26092049}; and at sizes
$6,7,8$ the first samples $16$ relations per size while the third samples
$64$. The sampled populations are therefore genuinely different, and that
independence is part of the evidentiary value.

\subsection{Exhaustive enumeration}

For $n$ labelled elements each program considers every reflexive relation by
choosing, independently, the $n-1$ possible off-diagonal entries of each row,
and then accepts exactly the transitive ones: whenever row $i$ contains $j$,
row $j$ must be a subset of row $i$. There are $2^{n(n-1)}$ candidates and
each labelled preorder is visited exactly once. No antisymmetry condition is
imposed, so preorders with genuine equivalence classes are included, and
isomorphism classes are not collapsed.

For every accepted relation \texttt{verify.py} independently computes the
maximal candidates from the ideal inequalities at total $n$, their exact
support counts by enumeration, and the alternating Boolean formula. It also
counts every interior-capacity case $\one+\one_T$ directly from the strict
filter inequalities, and compares the maximal candidates with all distinct
full-assignment vectors in the transport model.

\begin{center}
\begin{tabular}{@{}r r r r@{}}
\toprule
$n$ & Labelled preorders & Support identities & Interior counts\\
\midrule
0 & 1 & 1 & 1\\
1 & 1 & 2 & 2\\
2 & 4 & 16 & 16\\
3 & 29 & 232 & 232\\
4 & 355 & 5,680 & 5,680\\
5 & 6,942 & 222,144 & 222,144\\
\midrule
Total & 7,332 & 228,075 & 228,075\\
\bottomrule
\end{tabular}
\end{center}
All comparisons passed. The counts in the two right columns include zero
coefficients and empty support sets, not only successful nonempty supports,
and the interior column is checked at full exhaustive depth. The
assignment-model comparison was performed for each of the 7,332 preorders.

The supportwise verifier covers the same exhaustive range independently:
7,332 labelled preorders and 228,075 support identities, plus $64$ distinct
seeded preorders at each of sizes $6,7,8$, contributing another 28,672
identities for a total of 256,747. Its sampled relations are saved as row
bitmasks, so they can be recovered independently of the pseudorandom
generator. The transfer verifier's default exhaustive range is sizes
$0,\ldots,4$, that is $1+1+4+29+355=390$ relations; it accepts
\texttt{--max-n 5}.

\subsection{Additional cases and independent routes}

Using seed \texttt{20260920}, \texttt{verify.py} generates sixteen cases at
each of sizes six, seven and eight, alternating four construction modes:
sparse acyclic relations, dense acyclic relations, closures of arbitrary
directed relations, and partial orders on randomly chosen equivalence blocks.
In the recorded run all sixteen relations at each size were distinct. This
checks another 7,168 exact-support coefficients. The interior and assignment
checks were also run on the size-six and size-seven cases; they were omitted
at size eight in this test profile, which is stated openly rather than
implied.

For each of the five named examples a separate routine counts \emph{positive}
weighted multichains at $m=2,\ldots,n+2$, interpolates in the actual nodes
$\qint m$, evaluates at $-1/q$ and compares with the maximal-support
polynomial, without invoking the Boolean negative-evaluation formula. At
$q=1,2,3,5$ all twenty comparisons passed; for the source example the values
were
\[
 -18,\qquad -\frac{65}{32},\qquad-\frac{154}{243},\qquad-\frac{506}{3125}.
\]
This independent route is particularly useful for checking the $m-1$ chain
length and the $-2$ matrix exponent. The weighted lattice polynomial was also
checked in 156 evaluations on the five examples: twenty nonnegative capacity
vectors against direct lattice counts, zero capacities included, and 136
vectors in $\{-1,-2\}^E$ against the support--interior identity. Finally, the
nontransitive example of Section~\ref{sec:nontransitive} is explicitly
computed and checked to \emph{fail} the extension it is meant to refute.

The transfer verifier checks, over its exhaustive range: every coefficient of
the multivariate identity of Section~\ref{sec:proofB-support}; the full
transfer \eqref{eq:transfer} on deterministically generated signed integer
local polynomials, which is a sample of weights and not an exhaustive check
over all weights; and the negative evaluations as exact polynomials in $t$ for
$r=1,2,3,4$, the left side built from the finite product of
Lemma~\ref{lem:kernel} and the right side built independently from subsets of
colors and compositions, with the leading and smallest-degree claims of
Proposition~\ref{prop:degrees} checked at the same time. To avoid merely
comparing two rewritings of the same Gaussian expression, a further test
computes positive-length multichain weights by the direct dynamic program
\[
 V_1(a)=q^{|a|},\qquad V_{\ell+1}(a)=q^{|a|}\sum_{b\le a}V_\ell(b),
\]
and reconstructs the polynomial by exact Lagrange interpolation at $q=2,3$,
one further positive argument checking the interpolation before all four
negative arguments are compared against \eqref{eq:allnegative}.
Independent-capacity counts were checked by direct lattice enumeration for all
nonempty preorders of size at most three, all $u_i\in\{1,2\}$ and
$v\in\{0,1,2\}$, with reciprocity also checked at $v=1,2$.

\begin{table}[htbp]
\centering\small
\begin{tabular}{@{}lr@{}}
\toprule
Check family & Number passed\\
\midrule
Exhaustively enumerated labelled preorders, transfer verifier & 390\\
Full multivariate support identities in that range & 390\\
Signed local-polynomial transfer tests in that range & 390\\
Negative-evaluation polynomial identities, $r=1,\ldots,4$ & 1,560\\
Direct capacity counts and reciprocity evaluations & 1,250\\
Positive and negative interpolation comparisons & 410\\
Additional deterministic random preorders, sizes $5$--$8$ & 24\\
\bottomrule
\end{tabular}
\caption{Recorded output of \texttt{code/verify\_transfer.py} at its default
settings \texttt{--max-n 4}, \texttt{--max-r 4}, seed \texttt{20260920}. The
six worked examples also receive all structural and polynomial checks.}
\label{tab:checks}
\end{table}

The supportwise verifier adds checks that exist nowhere else. For every
preorder of size at most four, two triangular linear solves compute
$e_0^{\mathsf T}(ZD)^{-2}\mathbf j$ directly, at $q=2$, $q=3/2$ and with
independent coordinate weights $z_i=2,3,5,7$, for 1,170 matrix evaluations in
total; seven named examples are checked by interpolating the $q$-zeta
polynomial from positive multichain values at $q=2$, and the three-element
chain receives an eighth interpolation using the nonlinear height $h=|\cdot|^2$
with a degree bound of $9$. Above all, Lemma~\ref{lem:euler} is tested
\emph{standalone}, in isolation from the preorder application, on all 5,058
labelled bipartite graphs with $0\le|R|\le3$ and $0\le|C|\le4$ and on 300
further seeded graphs of sizes $(4,4)$, $(5,4)$, $(4,5)$, comparing three
independently computed counts: independent subsets of the private-element
graph, full-allocation degree vectors, and the integer vectors satisfying
\eqref{eq:interior-y}. All tests passed.

\subsection{Two finite algorithms, and what they cost}

Equation \eqref{eq:exact-support} gives a direct exact procedure for the whole
maximal-support polynomial: enumerate $B,S\subseteq E$, test
$\one_B+\one_S\in P_\tau$, and accumulate the signed sums. There are $4^n$
membership tests; using the allocation interpretation, each is a
bipartite-matching test on at most $2n$ demand copies and $n$ resources, and
sums exceeding $n$ can be rejected immediately. A basic augmenting-path
implementation gives an $O(4^nn^3)$ arithmetic-operation upper bound,
including all support counts. This is an exponential algorithm, not a claim of
polynomial-time enumeration in $n$.

Alternatively, for a fixed support $B$, enumerate the distinct degree vectors
of full allocations in $H_B$ by dynamic programming: begin with the zero
vector on $B$, and for each resource $r\in E\setminus B$ replace the current
set $D$ by
\[
 \bigcup_{s\in B:\,r\le_\tau s}\{d+e_s:d\in D\} .
\]
The cardinality of the final set is $N_\tau(B)$. The union is a \emph{set}
union: failing to deduplicate would count allocations rather than maximal
vectors. The shipped verifiers use a third engineering choice suited to small
exhaustive cases, encoding all ideals as a bitset and precomputing which
subset inequalities each candidate vector violates; this changes neither the
statement being checked nor its exact integer arithmetic.

If only the diagonal polynomial $H_\tau(u)$ is required, \eqref{eq:ternary}
uses $3^n$ vectors, each contributing $(-1)^{n+s}(1+u)^su^d$, and a direct
implementation takes $O(nj(\tau)3^n)$ elementary arithmetic operations, where
$j(\tau)$ is the number of ideals, with no construction of a large incidence
matrix. This bound ignores integer bit costs and is a transparent upper bound,
not a sharp complexity result. Against it, enumeration of all rank-$n$
candidate vectors requires at most $\binom{2n-1}{n-1}$ candidates when $n>0$,
which is often competitive. The value of the Boolean formula is therefore not
a universal speed claim; it is structural locality, the matroid
interpretation, and the ability to check two genuinely different computations
against one another.

\subsection{Running and reusing the programs}

From the extracted archive directory, the recorded profiles are
\begin{lstlisting}
python code/verify.py --exhaustive 5 --random-per-size 16 \
    --seed 20260920 --output data/verification.json
python code/verify_transfer.py --output data/verification_transfer.json
python code/verify_supportwise.py --output data/verification_supportwise.json
python -m unittest discover -s code -v
\end{lstlisting}
On PowerShell, put each command on one line, or replace the shell
continuation with PowerShell's continuation character. A smaller smoke test is
\begin{lstlisting}
python code/verify.py --exhaustive 3 --random-per-size 4 \
    --output data/smoke.json
\end{lstlisting}
Do not run Python with \texttt{-O}, since some test assertions use
\texttt{assert}.

The row-mask representation is explicit and shared by all three programs: bit
$j$ in \texttt{rows[i]} means $i\le_\tau j$, with zero-based labels. For
example, \texttt{(21, 26, 21, 8, 16)} encodes the five-element source example
in the order $a,b,c,d,e$; then
\texttt{boolean\_coefficients(rows)[24]} is $2$, the exact-support count for
$\{d,e\}$, and \texttt{direct\_qzeta\_negative(rows, 2)} is $-65/32$. Library
functions expose ideal generation, lattice points, maximal points, exact
support coefficients, cumulative support counts, direct interior counts, and
positive-chain interpolation. Invalid relations should be checked with
\texttt{is\_preorder}; transitive closure is available as a separate, explicit
operation, and a nontransitive relation must not be confused with its
transitive closure, which is never a silently applied repair.

The JSON records the actual relations used in the additional tests and a
SHA-256 digest of the deterministic exhaustive witness stream. That digest
helps detect accidental changes to the run; it is not a mathematical
certificate in the proof-assistant sense, and no checksum manifest of the
archive is shipped or intended.

The source paper already reports verification of its original, univariate
conjecture through size six~\cite[p.~15]{AC}. The role of these computations
is independent reproducibility and testing of the stronger support-by-support
identity, its two geometric engines, and the negative evaluations; it is not a
claim to exceed that original size bound. Neither the original checks nor
these establish an all-size theorem. That role belongs to the arguments in
Sections~\ref{sec:incidence}--\ref{sec:negative}.

\section{Proof audit and research status}
\label{sec:audit}

\subsection{Dependencies}

Two substantial results are imported: Postnikov's lattice-count formula
(Imported theorem~\ref{imp:lattice}) and Ehrhart--Macdonald reciprocity
\eqref{eq:ehrhart-reciprocity}. Engine~1 of Proof~A imports one more,
Postnikov's bipartite duality (Imported theorem~\ref{imp:duality}); Engine~2
and Proof~B do not. Hall's theorem is used by Proof~A and is proved in
Appendix~\ref{app:hall}; it is entirely absent from Proof~B. Everything else
is proved here: the Minkowski realization with independent capacities, the
filter identification of draconian vectors, the polynomial extension of
reciprocity, the incidence-matrix continuation and Boolean localization, the
support--interior identity, the private-element Euler identity, the residual
and privatization lemmas, the coefficient involution, the Gaussian endpoint
formula, and the positive local factors. The filter identification is included
with explicit credit to its appearance in \cite[Lemma~4.1]{AC}, and the
overlaps with \cite{PriorReciprocity} are listed in
Section~\ref{sec:nonduplication} and flagged again at
Remarks~\ref{rem:prior-enumerator} and~\ref{rem:prior-reciprocity}. No theorem
asserted by another report in the supplied manifest is used as a proof source.

\subsection{Normalization errors the proofs are sensitive to}

The following points make the delicate conventions explicit. They merge the
audits of the three source drafts; where two drafts recorded the same point
differently, both formulations are kept.

\begin{enumerate}[label=\arabic*.]
\item \textbf{Lower covers and rank are different.} Covers are $a-e_i$ for
positive coordinates; the exponent in the conjecture is the support size, not
the common maximal rank $n$.
\item \textbf{The matrix is $ZD$, and the exponent is $-2$.} Its column
weighting gives the sum of the heights of the chain entries. There are $m-1$
entries, so the negative evaluation uses $A^{-2}$, not $A^{-1}$. The
independent multichain-interpolation tests were designed specifically to
detect an error here.
\item \textbf{Repeated height eigenvalues cause no Jordan blocks.} The rank
(or height) filtration proves annihilation by the \emph{squarefree}
polynomial $\prod_j(X-q^{g_j})$ over $\Q(q)$. The argument does not rely on
upper triangularity alone.
\item \textbf{Boolean localization localizes only the middle vector.} The
M\"obius factor $\mu(\zero,a)$ forces $a$ to be Boolean, but $a+\one_S$ may
have entries equal to $2$ and must not be treated as Boolean; the signed count
does not discard those terms.
\item \textbf{The height statistic is $h(\one_{\supp(b)})$, not $h(b)$.} The
nonlinear-height example on the three-element chain distinguishes the two.
\item \textbf{The weighted count uses the dual preorder.} Its source
neighbourhoods are $D_\tau(i)$, not $U_\tau(i)$. This orientation is checked
twice, in Proposition~\ref{prop:count} and again in the support-restricted
graph of Section~\ref{sec:interior}.
\item \textbf{The zero-weight universal vertex remains present.} In the
untrimmed Postnikov formula its factor is $f_{a_0}(0+1)=1$, not
$\binom{a_0-1}{a_0}$; the new $0$-coordinate is included when applying the
theorem, so the draconian total is $n$ and not $n-1$. Deleting that vertex
would change both the total and the formula.
\item \textbf{Interior strictness is integral strictness.} After subtracting
$\one$, the capacity $\one+\one_T$ produces $y(I)\le|I\cap T|-1$, not
$y(I)\le|I\cap T|$; likewise in Engine~2 the shift by $\one_R$ removes the
private cube and leaves $y(J)\le|N(J)|-1$.
\item \textbf{The support-restricted graph need not be connected initially.}
The isolated-vertex case has both counts zero; in the remaining case the
universal right neighbour makes the graph connected, meeting the hypothesis of
the imported duality. Correspondingly, in Engine~2 no empty simplex is ever
passed to the nonempty-simplex formula: a resource with no nonprivate
neighbour is handled before any polytope is formed.
\item \textbf{Degree vectors, not allocations.} The objects in bijection with
maximal points are distinct full-allocation degree vectors. Counting maps or
matchings with multiplicity would give the wrong theorem, which is also why
the dynamic program of Section~\ref{sec:computation} must deduplicate.
\item \textbf{The multivariate step is not a weighted $q$-Ehrhart
assertion.} It is the explicit two-value identity of
Lemma~\ref{lem:interpolation}, followed by ordinary univariate Ehrhart
reciprocity along each fixed capacity ray.
\item \textbf{Capacities are independent variables before any negative
specialization.} In Proof~B reciprocity is established in all $n+1$ variables
first, and only then is $v$ set to $0$; there is no polytope of capacity $-1$.
\item \textbf{Preorders, not just partial orders.} Privatization uses
reflexivity and transitivity and never antisymmetry; genuine equivalence
classes are allowed and retain their individual coordinates. Boundary cases
with empty $E$, empty residual resource set and empty target sets are treated
explicitly. The exhaustive verifiers include relations with genuine cycles.
\end{enumerate}

\subsection{What is established, and what is not}

The arguments supplied here prove Theorems~\ref{thm:main},
\ref{thm:allnegative} and~\ref{thm:height-intro} from the stated classical
results, with no unproved additional combinatorial assertion and no hidden
lemma. In that mathematical sense this is a complete proposed solution of the
chosen conjecture, and its validity is not inferred from the finite
computations.

The following separate questions are \emph{not} answered by the computations
or by this presentation: independent expert verification of every step;
priority relative to unpublished or newly posted work; and a formal proof
checked by a small trusted kernel. No such verification or priority
certificate is claimed. This is an AI-assisted research draft, not
independently refereed, not formalized in a proof assistant, and not
represented as accepted or published mathematics. The literature checks were
performed on 20 September 2026 and are documented in
Appendix~\ref{app:ledger} and \texttt{notes/source\_audit.md}. The target is
identified by both conjecture number and displayed formula, so a future change
of source numbering need not make the statement ambiguous. It should be added
that three separate drafts reaching the same two statements by three different
routes is evidence of internal consistency, not of correctness; a shared
misreading of the source's conventions would not have been detected by that
agreement alone, which is why those conventions are restated explicitly in
Section~\ref{sec:problem} and audited above.

\subsection{Provenance of the source discrepancy}
\label{sec:discrepancy}

All three merged drafts recorded the same anomaly about the target's version
metadata, each slightly differently, and the union is retained. The arXiv
submission record displayed only v1, submitted 26 May 2026; the PDF's author
date is 27 May 2026; and the retrieved HTML bears an internal date in
August 2026 --- one draft recording it as 24 August 2026. The three drafts
therefore pin the target by arXiv identifier, section, equation number and
displayed formula rather than by any of those dates, and the PDF's conjecture
page was inspected visually as well as through parsed text. The relevant
$q$-zeta statement agrees between all representations consulted.

Two further provenance points are kept verbatim from their original drafts.
First, only the manifest, and not the full text of the earlier packages, was
available when nonduplication was assessed. Second, the manifest's relevant
entry is preserved word for word, with its SHA-256 digest and its source line
numbers 960--969, in \texttt{notes/manifest\_provenance.txt}; that file is the
precise record of the nonduplication boundary, and it is the only place where
the earlier entry is reproduced exactly.

\section{Conclusion}

The negative $q$-zeta evaluation is controlled by a small-coordinate part of
the lattice-point poset, but its positive interpretation is global: maximal
points of total $n$ with a specified support. Proof~A links the two views
geometrically. In Engine~1 the support-restricted graph $G_T$ does the work:
its left draconian sequences encode maximal points allowed to use $T$, its
right sequences encode interior points at capacities $\one+\one_T$, classical
reciprocity converts the latter to negative polynomial values, and a two-entry
interpolation table produces the same local weights as incidence inversion. In
Engine~2 the work is done instead by privatizing the residual targets and by
an Euler identity valid for every bipartite graph, which trades length for a
smaller imported surface. Proof~B links the two views algebraically, with no
geometry beyond the one reciprocity statement: an auxiliary capacity is kept
alive, then set to zero, and an involution on local power series moves the
whole computation onto the maximal layer --- which is also why it, and it
alone, reaches every $\qint{-r}$.

Together these yield the full Athanasiadis--Chapoton $q$-reciprocity formula,
the stronger exact-support identity \eqref{eq:exact-support}, its
arbitrary-height version, and positive integral formulas at every negative
$q$-integer. The three obstructions of Section~\ref{sec:limits} show that the
cancellation is not an automatic property of pure coordinate downsets, of
arbitrary assignment models, or of arbitrary capacities. It is the transitive
ideal--filter structure of a preorder, with unit capacities, that makes the
different descriptions agree.

\appendix

\section{The finite Hall theorem used in Proof A}
\label{app:hall}

We recall the standard finite theorem~\cite{Hall}, with its short inductive
proof, in order to fix the exact matching convention. Let $L,R$ be finite
disjoint sets and let $N(A)\subseteq R$ be the neighbours of $A\subseteq L$.
There is an injective assignment $f:L\to R$ with $f(i)\in N(\{i\})$ if and
only if
\[
 |A|\le|N(A)|\qquad(A\subseteq L).
\]
Necessity is immediate. For sufficiency, induct on $|L|$; the empty case is
trivial. If some nonempty proper $A\subset L$ has $|N(A)|=|A|$, match $A$ into
$N(A)$ by induction; for $B\subseteq L\setminus A$,
\[
 |N(B)\setminus N(A)|=|N(A\cup B)|-|N(A)|\ge|B| ,
\]
so induction matches $L\setminus A$ into $R\setminus N(A)$ and the two
matchings combine. If no such nonempty proper $A$ exists, choose $i\in L$ and
a neighbour $r\in N(\{i\})$; after deleting $i$ and $r$, every nonempty
$B\subseteq L\setminus\{i\}$ still has at least $|B|$ neighbours, since
originally it had at least $|B|+1$. Apply induction and add the pair $(i,r)$.
The singleton case is included directly by the existence of a neighbour.

Every use with integer multiplicities replaces a demand $i$ of multiplicity
$a_i$ by $a_i$ distinct copies. For any subset of copies, its cardinality is
at most the total multiplicity of its set of types, while its neighbourhood is
the neighbourhood of that set of types. This licenses every passage from
vector inequalities to an ordinary finite matching in this document.

\section{A compact verification specification}
\label{app:spec}

A reimplementation need not reproduce the Python organization. The essential
checks are the following, all over the integers except the last.

\subsection{Support, Boolean sums, and interior counts}

For each enumerated preorder let $\mathcal J$ be its ideals and $\mathcal F$
its filters. For every $B,T\subseteq E$ compare
\begin{align*}
 h(B)&=\#\{a\in\N^E:|a|=n,\ \supp(a)=B,
                         \ a(I)\le|I|\ (I\in\mathcal J)\},\\
 h(B)&\stackrel{?}{=}(-1)^{n-|B|}
      \sum_{S\subseteq E}(-1)^{|S|}
      \one\bigl[|B\cap I|+|S\cap I|\le|I|\ (I\in\mathcal J)\bigr],\\
 \sum_{B\subseteq T}h(B)&\stackrel{?}{=}
      \#\{y\in\N^E:y(I)\le|I\cap T|-1\ (\emptyset\ne I\in\mathcal F)\}.
\end{align*}
The full-set bound makes the last enumeration finite, with $|y|\le|T|-1$; when
$T=\emptyset$ and $n>0$ its count is zero. The additional transport check
compares the set of rank-$n$ points above with
\[
 \left\{\sum_{i\in E}e_{f(i)}:f:E\to E,\ i\le_\tau f(i)\text{ for every }i
 \right\},
\]
different assignments giving the same vector being counted only once.

\subsection{The independent \texorpdfstring{$q$}{q}-check}

Enumerate all points $a$ of total at most $n$ satisfying the ideal bounds. At
a fixed positive integer $q$, initialize $v_1(a)=q^{|a|}$ and iterate
$v_{\ell+1}(b)=q^{|b|}\sum_{a\le b}v_\ell(a)$; then $\sum_av_{m-1}(a)$ is the
positive value at $\qint m$. Interpolate a polynomial of degree at most $n$
through the $n+1$ nodes $m=2,\ldots,n+2$ using rational arithmetic and compare
its value at $-1/q$ with $(-1)^n\sum_Bh(B)q^{-|B|}$. At $q=1$ use
$\qint m=m$. The interpolation must use exact rational arithmetic even at
negative $q$-integers such as $\qint{-2}=-3/4$ at $q=2$.

\subsection{Capacities, transfer, and negative evaluations}

A relation is stored as $n$ integer bit masks, bit $j$ of row $i$ being one
exactly when $i\le_\tau j$; transitivity is verified by requiring every bit of
row $j$ to be present in row $i$ whenever bit $j$ is present in row $i$.
Ideals are generated by testing downward closure of every subset. This
representation retains labels and does not collapse equivalence classes.

The capacity test implements \eqref{eq:E} with the falling-factorial
definition of the binomial coefficient, including negative upper arguments.
The transfer test applies the triangular integral rule \eqref{eq:Tcoeff} to
signed integer local polynomials. For a univariate polynomial the JSON
coefficient array is in ascending powers, so $[0,2,9,6,1]$ means
$2t+9t^2+6t^3+t^4$. Each example records the relation rows, the total
lattice-point count, the number and full list of maximal points, the support
polynomial and all four signed negative evaluations.

\section{Compact proof outlines for independent checking}
\label{app:outline}

Proof~A compresses to the four-line chain displayed in
Section~\ref{sec:completion}, each line labelled with the ingredient it uses.
Proof~B compresses to the following five identities, starting from the count
polynomial \eqref{eq:E}:
\begin{align*}
 \Ecal_\sigma(-u,-v)
  &=(-1)^n\Ecal_\sigma(u-\one,v-1),\\
 \sum_{a\in P_\tau}\prod_i\binom{a_i-u_i-1}{a_i}
  &=(-1)^n\sum_{b\in \Max(P_\tau)}\prod_i\binom{b_i+u_i-2}{b_i},\\
 \sum_{a\in P_\tau}\prod_i[z^{a_i}]f_i(z)
  &=(-1)^n\sum_{b\in \Max(P_\tau)}\prod_i[z^{b_i}](\Tr f_i)(z),\\
 q^a\qbinom{a-r-2}{a}
  &=[z^a]\prod_{j=0}^r(1-q^{-j}z),\\
 \Tr\!\left(\prod_{j=0}^r(1-t^jz)\right)
  &=\prod_{j=1}^r\left(1+\frac{t^jz}{1-z}\right).
\end{align*}
The first line is interior-point translation plus classical reciprocity; the
second is $v=0$ with the dual preorder; the third is basis interpolation; the
fourth is a finite Gaussian identity; the fifth is one rational substitution.
For $r=1$ the last expression has coefficient $t$ in every positive degree,
which is exactly the support statistic in the conjecture.

Engine~2 of Proof~A compresses to: $\eta_\tau(B)$ is an alternating sum over
the independent sets of a private-element bipartite graph
(Lemmas~\ref{lem:residual} and~\ref{lem:private}); that alternating sum is
$L_K(-1)$ for the polytope $K=[0,1]^R+\sum_j\Dz{A_j}$; Ehrhart reciprocity
turns it into an interior count; a second application of Postnikov's theorem
turns that into $\#(\sum_i\Delta_{N(i)}\cap\Z^{C})$; and Hall's theorem
identifies those lattice points with full-allocation degree vectors, which are
in bijection with the maximal points of support $B$.

\section{Source and version ledger}
\label{app:ledger}

The mathematical sources and the claimed additions are separated as follows.
This ledger merges the three ledgers of the drafts named in
Section~\ref{sec:provenance}; where they disagreed, the more cautious
statement is kept.

\begin{description}[style=nextline,leftmargin=1.2em,labelwidth=0pt,labelindent=0pt,font=\normalfont\bfseries]
\item[Athanasiadis--Chapoton, arXiv:2605.26916v1.]
Section~4.4, Conjecture~4.9 and equation~(21), on printed page~14 of the
retrieved PDF, supply the target; Section~3.1 supplies the five-element
example and its defining inequalities; Theorem~2.1 and Lemma~4.1 are the
source's own statement of Postnikov's formula and of the draconian/filter
identification. The source's printed page~15 reports computational
verification through size six and the value $-(2q+1)(q^2+4q+1)/q^5$. The
version-date anomaly is recorded in Section~\ref{sec:discrepancy}. The PDF's
conjecture page was inspected visually as well as through parsed text.
\item[Chapoton, arXiv:2402.11979v2, revised 29 January 2025.]
Supplies the $q$-zeta concept, its conventions, and the general
height-weighted definition; Appendix~C develops the $q$-incidence algebra of
which Section~\ref{sec:incidence} is a written-out specialization. It is not
claimed that interpreting $q$-zeta values through incidence powers is new. No
general weighted reciprocity theorem is assumed from this paper: the negative
specializations used here are derived independently, twice
(Remark~\ref{rem:two-constructions}).
\item[Postnikov, arXiv:math/0507163v1; IMRN 2009(6), 1026--1106.]
Theorem~11.3, printed page~28 of the arXiv version, supplies the lattice-point
formula, with its untrimmed full-simplex shift retained exactly;
Definition~9.2, Lemma~11.7 and Corollary~11.8 supply the bipartite duality
used only by Engine~1. The pages containing both were checked in the PDF page
images. The passage of that paper explicitly called an ``alternative
semiproof'' is not what is imported; the stated theorem and its main proof
are the dependency. These are explicit imports, not part of any novelty claim.
\item[Ehrhart--Macdonald reciprocity.]
Used only for full-dimensional integral polytopes with all capacities positive
integers, univariately along $s\mapsto s\mathbf c$. Three references are
retained: Macdonald~\cite{Macdonald} as the historical source,
Beck--Robins~\cite[Chapter~4]{BeckRobins} as the textbook statement, and
Beck--Develin~\cite{BeckDevelin} as a primary-source proof of the rational
cone reciprocity from which the polytope case follows. The Macdonald publisher
page was not accessible through the browsing tool used, and that bibliographic
entry was cross-checked indirectly in research-paper references; no claim is
made that the original Macdonald paper was read in full.
\item[Hall, JLMS 10 (1935), 26--30.]
The finite marriage theorem, also proved in Appendix~\ref{app:hall}.
\item[Dai--Hou--Liu--Thawinrak--Wang, arXiv:2608.16037v2, 27 August 2026.]
Checked as related recent work. Its preorder application concerns duality for
the support enumerator of \emph{all} lattice points, which is a different
enumerator from the maximal-support polynomial here; see the $V$ versus $V^*$
example of Section~\ref{sec:examples}. It is not a proof dependency. A text
search for ``zeta'' in the retrieved HTML returned no matches; that is a
search observation, not a certificate of absence of related mathematics.
\item[\nolinkurl{enumerative-combinatorics/preorder-polytope-reciprocity}.]
The earlier package \cite{PriorReciprocity} in this collection, cited in the
three places listed in Section~\ref{sec:nonduplication}. It is credited as
prior work, not imported as a lemma, and not audited here.
\item[User-supplied \texttt{manifest.tex}, 20 September 2026.]
Used only to delimit the target and to assess nonduplication, never as a proof
source. Only the manifest, and not the full text of the earlier packages, was
available for that comparison. The relevant entry is preserved verbatim with
its digest and line numbers in \texttt{notes/manifest\_provenance.txt}.
\item[Status searches, 20 September 2026.]
Targeted searches combined the source title or arXiv identifier with
``Conjecture 4.9'', ``$q$-zeta'', ``reciprocity'' and ``proof''; the source
manuscript, Chapoton's paper and the cross-polytope paper were inspected. The
checked source still states the target as a conjecture, and no separate
resolution was located. Some results were irrelevant or aggregated. This is
evidence about the sources checked, not a universal priority claim: an
unpublished proof, an unindexed announcement, or the same result in other
terminology could have been missed.
\end{description}

\section{Suggested manifest entry}
\label{app:manifest}

The included file \texttt{manifest\_entry.tex} contains a ready-to-adapt
entry. Its scope is deliberately limited to the full $q$-refinement of
Conjecture~4.9, the stronger support identity, the evaluations at every
negative $q$-integer and the arbitrary-height variant. It does not repeat the
earlier packages' other claimed conjecture resolutions, and its wording
retains the research-draft status instead of representing this investigation
as independently refereed. The three drafts merged here proposed three
different directory paths, none of which matched the actual location of any of
them; the entry now names one correct path.

\begin{thebibliography}{9}
\bibitem{AC}
C.~A. Athanasiadis and F. Chapoton,
\emph{Polytopes and posets associated to preorders},
arXiv:2605.26916v1, 2026.
\url{https://arxiv.org/abs/2605.26916v1}.
Target: Section~4.4, Conjecture~4.9, equation~(21), printed page~14;
also Theorem~2.1, Lemma~4.1 and Section~3.1.

\bibitem{Chapoton}
F. Chapoton,
\emph{On a $q$-analogue of the Zeta polynomial of posets},
arXiv:2402.11979v2, revised 29 January 2025.
\url{https://arxiv.org/abs/2402.11979v2}.
The $q$-incidence algebra is developed in Appendix~C.

\bibitem{Postnikov}
A. Postnikov,
\emph{Permutohedra, associahedra, and beyond},
International Mathematics Research Notices \textbf{2009}, no.~6,
1026--1106. DOI:~10.1093/imrn/rnn153.
\url{https://arxiv.org/abs/math/0507163}.
Theorem and lemma numbering in this draft follows the arXiv version:
Definition~9.2, Theorem~11.3, Lemma~11.7 and Corollary~11.8.

\bibitem{Macdonald}
I.~G. Macdonald,
\emph{Polynomials associated with finite cell-complexes},
Journal of the London Mathematical Society (2) \textbf{4} (1971), 181--192.

\bibitem{BeckRobins}
M. Beck and S. Robins,
\emph{Computing the Continuous Discretely: Integer-Point Enumeration in
Polyhedra}, second edition, Springer, 2015, Chapter~4.
DOI:~10.1007/978-1-4939-2969-6.

\bibitem{BeckDevelin}
M. Beck and M. Develin,
\emph{On Stanley's reciprocity theorem for rational cones},
arXiv:\allowbreak math/\allowbreak 0409562v3, revised 4 August 2005,
available at \url{https://arxiv.org/abs/math/0409562v3}. Only the
lattice-polytope consequence is used here, with full dimensionality checked
explicitly at each application.

\bibitem{Hall}
P. Hall,
\emph{On representatives of subsets},
Journal of the London Mathematical Society \textbf{10} (1935), 26--30.
DOI:~10.1112/jlms/s1-10.37.26.

\bibitem{DaiEtAl}
Z. Dai, Q. Hou, Z. Liu, W. Thawinrak and H. Wang,
\emph{Counting Lattice Points in Minkowski Sums of Cross Polytopes},
arXiv:2608.16037v2, 27 August 2026.
\url{https://arxiv.org/abs/2608.16037v2}.

\bibitem{PriorReciprocity}
\emph{Reciprocity and Matrix Duality for Preorder Polytopes},
research manuscript, 19 September 2026, in this collection under
\nolinkurl{docs/reports/enumerative-combinatorics/} in the directory
\nolinkurl{preorder-polytope-reciprocity}. It proves Conjectures~7.2, 4.2
and 4.4 of \cite{AC} and the ordinary $q=1$ case of Conjecture~4.9, together
with multivariate reciprocity for independently weighted capacities. Cited as
prior work in this collection; not imported as a lemma and not audited here.

\bibitem{Manifest}
\emph{Manifest of mathematical research packages},
user-supplied \texttt{manifest.tex}, received 20 September 2026.
Relevant entry: ``Reciprocity and Matrix Duality for Preorder Polytopes,''
source lines 960--969. The manifest is used only to delimit the present task,
not as a proof source. A verbatim entry with its digest accompanies this
archive in \texttt{notes/manifest\_provenance.txt}.
\end{thebibliography}

\end{document}
